\documentclass[../main.tex]{subfiles}
\begin{document}
%==========================================TIÊU ĐỀ TẬP========================================
\begin{table}[h]
  \begin{adjustwidth}{-1cm}{}
    \begin{tabular}{>{\centering\arraybackslash}p{6cm}|>{\raggedright\arraybackslash}p{12.5cm}}
      \multirow{5}{6cm}
      % Cột bên trái: ÉPISODE và số tập
      {\\[-40pt]
        \flaregothic\fontsize{20pt}{20pt}\selectfont \centering
        {\contour{errorfive}{\textcolor{black}{ÉPISODE}}}\color{black}\\
        \burbank\fontsize{72pt}{72pt}\selectfont
        \includegraphics[height=3.12359cm]{../../../../Image/Book?/number/num26.png}
        \\[18pt]
        % \flaregothic\fontsize{14pt}{14pt}\selectfont
        % \color{epattr}BIENVENUE, \uline{\color{Banpire} Mel} !
        % \flaregothic\fontsize{18pt}{18pt}\selectfont
        % \color{epattr}ÉPISODE PERDU
        \color{black}
      }
       &
      % Dòng thứ nhất: tên tập tiếng Nhật
      {\fotskip\fontsize{18pt}{18pt}\selectfont \ruby{日}{に}\ruby{常}{じょう}の\ruby{学}{がっ}\ruby{校}{こう}\ruby{生}{せい}\ruby{活}{かつ}}
      \\[6pt]
      % Dòng thứ hai: tên tập tiếng Anh
       & {\cafeteria\fontsize{18pt}{18pt}\selectfont Daily School Life}                                                                                                                                                    \\[4pt]
      % Dòng thứ ba: tên tập tiếng Việt
       & {\vnmsans\fontsize{18pt}{18pt}\selectfont Cuộc sống học đường thường nhật}                                                                                                                                        \\[8pt]
      % Dòng thứ tư: nhân vật xuất hiện
       & {\caviar{\fontsize{14pt}{14pt}\selectfont Nhân vật: Toàn bộ các nhân vật xuất hiện trong \hll\ ERROR, với các nhân vật mới (\emp{choco} \emp{azki} \emp{haato} \emp{aqua} \emp{shion} \emp{ayame} \emp{watame})}}
      % \\[6pt]
      % Dòng cuối cùng: Thời gian diễn ra của tập (thường sẽ là ngày phát hành)
      % & {\continuum\fontsize{16pt}{16pt}\selectfont Ngày phát hành: 11/06/2022}\\[8pt]
    \end{tabular}
  \end{adjustwidth}
\end{table}
%===========================================NỘI DUNG==========================================
\color{black}

\begin{center}
  \pov{Hikawa Kuon}
  \uline{Nhà Hikawa (tạm) - thị trấn Aogami.\\Ngày hôm sau\\Năm giờ sáng.}
\end{center}

\dots Chậc.

Tiếng chuông báo thức quen thuộc vang lên từ điện thoại của tôi, vẫn là cái bản nhạc ``ám ảnh'' đó mà tôi đã cài suốt nhiều năm nay. Chỉ cần nghe vài nốt đầu là tôi tỉnh ngay lập tức, như thể não tôi đã được lập trình để bật dậy ngay khi giai điệu ấy vang lên.

Tôi dụi mắt, ngồi dậy, quay sang nhìn hai chị em đang cuộn tròn trong chăn. Cả hai thở đều đều, gương mặt bình thản đến mức khiến tôi tự hỏi không biết họ đang ngủ thật hay đang cố tình ngủ say để trốn tiếng chuông hay không.

Tôi hắng giọng: ``\vie{Này, dậy đi. Muộn giờ học bây giờ.}''

\dots Không nhúc nhích.

Suzuki chỉ xoay người một chút, ôm chặt cái gối hơn.

Kaigou-chan thì mặt chôn trong chăn như người đã thoát ly khỏi thế giới này.

Tôi thở dài, đứng dậy, bắt đầu chuẩn bị đồng phục và cặp sách của mình. Thỉnh thoảng liếc sang hai người vẫn ngủ không biết trời trăng mây gió ra sao.

Khi tôi gọi lần thứ hai, họ vẫn chẳng có phản ứng gì ngoài tiếng thở nhè nhẹ.

``\dots Cạn lời với hai người luôn.''

Khi đã chuẩn bị xong, tôi bỏ cuộc luôn việc kéo họ dậy, xuống bếp làm bữa sáng. Tôi chẳng kì vọng gì, nhưng vẫn theo thói quen mà làm luôn hai hộp cơm để mang đi. Bên trong cũng đơn giản thôi: trứng cuộn, xúc xích, cơm rắc mè, một chút rau luộc, và đâu đó đồ còn sót từ bữa cơm tối ngày hôm qua.

Lên lầu gọi lần thứ ba. ``\textbf{Dậy ngay!}''

Nhưng đáp lại tiếng hét của tôi chỉ là những câu nói trong mơ.

Kaigou-chan mơ màng lẩm bẩm: ``\dots Mười phút nữa\dots chỉ mười phút\dots''

Suzuki thì ú ớ: ``\dots Đừng\dots có ăn\dots bánh mì của mình chứ\dots''

Tôi khoanh tay, nhìn hai cái kén ngủ nằm bất động, cảm giác như mình vừa bị một bức tường vô hình đẩy ngược trở lại. ``\vie{\dots Thôi, kệ hai người luôn đấy.}''

Tôi học kiểu đó từ bố tôi ở thế giới gốc. Ông luôn bảo: ``\vie{Gọi không dậy thì cứ để đó, đói rồi tự biết đường mà sửa.}'' Lời dạy nghe hơi phũ, nhưng hiệu quả đến mức tôi rèn luyện được tính đúng giờ này từ lúc nào không biết.

Tôi xuống nhà, ăn sáng một mình. Trước khi đi, tôi đặt vào đó hai tờ giấy nhắn nhỏ trên hai chiếc hộp cơm.

Trên hộp của cô chị, tôi viết:

\txtbx{``Có muộn đến mấy cũng phải ăn, kẻo học giữa chừng đói lả. Nhớ khóa cửa hộ tôi.''}

Còn trên hộp của em gái cô ấy, tôi để lại:

\txtbx{``Nếu em để đói rồi xỉu thì đừng đổ thừa. Lần sau cố mà dậy sớm.''}

Đặt hai chiếc hộp ngay ngắn lên bàn, tôi khoác cặp lên vai và bước ra ngoài.

\ct

Đi tới trường một mình\dots\ cảm giác tôi đã quá quen từ thời trung học ở thế giới gốc. Khoảng thời gian đó, tôi luôn đi sớm nhất nhà, lang thang qua từng con phố mờ sương, thinh lặng đến mức tiếng bước chân mình nghe rõ mồn một. Thật kỳ lạ, dù thế giới này khác thế giới kia, không khí sáng sớm vẫn mang một màu rất giống nhau.

Giọng Game vang lên khe khẽ trong chiếc đồng hồ: ``\eng{Hoài niệm sao, Kuon-user?}''

Tôi rướn mày: ``\eng{Có chút. Trước khi vào \hll, mỗi ngày tôi đều như thế này.}''

``\eng{Bây giờ thì sao?}''

``\eng{Có một chút khác biệt. Bây giờ tôi không còn một mình nữa.}'' Tôi bật cười nhỏ. ``\eng{Chỉ là hai người kia còn ngủ chết dí thôi.}''

Game cười nhạt trong âm sắc trầm đặc: ``\eng{Chẳng trách cậu được. Họ dậy muộn thì tự chịu thôi.}''

Đến gần khu phố dẫn vào trường, tôi gặp Yae-san và Hotaru-san đang đi cùng nhau. Tôi gật đầu chào: ``Chào buổi sáng, hai người.''

Yae-san cười nhẹ, hai khuỷu tay xoay đều như đang tập thể dục: ``Kuon-san dậy sớm thật đó.''

Hotaru-san thì giơ tay chào, nhìn có vẻ vẫn còn ngái ngủ. Tôi hỏi trêu: ``Đêm qua ngủ muộn hả?''

``\dots Một chút,'' cô bạn da ngăm đáp.

``Như thế là không được đâu đó, Hotaru-san. Con gái con đứa không nên thức khuya chứ,'' cô gái tóc xám nhắc nhở, đôi mày nhíu lại.

Tôi tiếp tục đi về phía trước, để hai người họ đối đáp nhau.

Vừa mới tiếp tục bước được vài bước, thì bất chợt\dots

``\textbf{Xin lỗi, xin lỗi! Cho tôi đi qua!!}''

Một bóng hồng-xanh vụt ngang qua trước mặt tôi, nhanh đến mức gió tạt vào má. Chẳng mấy chốc\dots

``Kya!?''

*RẦM*

Cô gái ngã sấp mặt giữa đường, đôi bím tóc tết hồng-xanh đung đưa đáng thương. Tôi đứng chết lặng nửa giây trước khi bóp trán ngán ngẩm. ``\dots Tới cả thế giới này mà cô ấy vẫn còn giữ tính hậu đậu được\dots''

Tôi bước tới, đỡ cô dậy, phủi bụi trên đồng phục cô theo phản xạ tổng quản hololive.

Cô gái đỏ mặt từ tai tới cổ, lắp bắp: ``X-xin lỗi!! Mình vội quá!!''

Rồi chạy biến như gió, để lại tôi và một khoảng trời im lặng kỳ dị.

Tôi thở dài: ``Tính cách của cô ấy cũng chẳng khác Aku-tan là bao.''

``\eng{Người cậu quen sao, Kuon-user?}'' giọng nói trong chiếc đồng hò đeo tay chợt lên tiếng.

``\eng{\dots Không hẳn. Nhưng mà cái kiểu ăn nói ấy, tính cách hậu đậu kia\dots\ không trật đi đâu được.}''

Tôi có cảm tưởng rằng tôi sẽ gặp lại cô ấy khá sớm đấy. Linh tính tổng quản của tôi mách bảo vậy.

\ct

Tôi tiếp tục đi, qua cổng trường vừa lúc nắng sáng bắt đầu chiếu qua hàng cây.

Và tại nơi này, chuyện rắc rối đầu tiên của buổi sáng đang chờ tôi phía trước.

Phải rồi. Đã là đời học đường thì kiểu gì cũng phải có đám bắt nạt.
Tôi còn lạ gì cái mô-típ này nữa — tuy trường tôi học ở thế giới gốc chẳng bao giờ có thể loại này, nhưng văn hoá đại chúng thì tôi đọc đủ rồi, báo chí cũng đọc không thiếu.

Vậy nên khi tôi vừa bước qua cổng trường và thấy một nhóm con trai đứng chắn ngang lối đi, mặt thằng nào cũng hằm hằm như vừa bị tước mất bữa sáng\dots\ tôi chỉ biết thở dài: ``Đúng như kỳ vọng của tôi thật.''

Thằng cầm đầu khoanh tay, nhìn tôi từ trên xuống, mặt hằm hằm: ``Mày là Hikawa Kuon đúng không?''

``Ừ. Có chuyện gì à, các cậu?''

Câu ``các cậu'' vừa rời khỏi miệng là tôi thấy thái độ cả nhóm co giật. Với tôi thì là lịch sự, nhưng với chúng nó, hẳn là nghe như bị coi thường.

Một thằng khác gằn giọng: ``Tao hỏi thật\dots\ mày nghĩ sao khi được xếp vào cái lớp toàn con gái đó?''

À, hiểu rồi. Sự ghen tị.

Tôi đáp tỉnh bơ, chẳng cần diễn: ``Thật ra thì\dots\ việc trường xếp tôi vào lớp full con gái cũng khiến tôi bất ngờ lắm chứ.''

Tôi nhún vai, tiếp tục: ``Nhưng mà nói thật, nếu mấy cậu có vào lớp đó thì cũng chưa chắc họ thích mấy cậu đâu.''

Cả nhóm đồng loạt quay sang nhìn tôi như thể tôi vừa cầm dao đâm vào lòng tự trọng của chúng.

Tôi tiếp tục, giọng đều đều: ``Vì, thành thật mà nói, \emph{bản thân mấy cậu còn chẳng ra gì} nữa là. Họ chẳng đếm xỉa tới những tên thích thể hiện như mấy cậu đâu.''

``Mày--!''

Một thằng lao lên nhưng bị thằng cầm đầu chặn lại.

Tên trùm tiến gần, nắm lấy cổ áo tôi: ``Mày nói lại xem?''

Tôi nhìn xuống tay nó, rồi thở dài, nhìn lên mặt nó: ``Bài này cũ quá rồi, cậu biết không? Sao không cho tôi cái gì đó mơi mới xem?''

Hắn nghiến răng, giơ nắm đấm lên. Nhanh đấy, nhưng không nhanh bằng những lần tôi phải né mấy quả đạn pháo, mấy cú đấm trời giáng, đất đá bê tông từ những lần phá hủy tòa nhà, hay cả đĩa sushi biết bay của các nàng hề tôi dẫn dắt.

Tôi chỉ nghiêng đầu nhẹ sang một bên\dots

*BỐP!*

\dots là nắm đấm của hắn đi thẳng vào mặt đứa đồng bọn phía sau. Thằng đó lảo đảo, chửi tục: ``\textbf{Chó thật, mày đấm nhầm rồi, khốn nạn!}''

Tên trùm mất đà đá ngang, tôi chỉ hơi nhún chân nhảy lên. Cú đá bay thẳng vào bụng chính đứa đang định nhào tới phụ hắn. Thằng này ôm bụng gục xuống, rên lên như em bé: ``Ơi\dots\ đau quá\dots''

Tôi đáp đất nhẹ nhàng, cười khẩy: ``Ôi, hai cú đấm đá mà mấy cậu mất hai tên hậu cần rồi à?''

Cả bọn đồng loạt nhìn tôi như thể tôi là sinh vật ngoài hành tinh.

Một đứa phía sau quát: ``\textbf{Để tao lo thằng này!}'' rồi lao tới từ góc chết phía sau lưng tôi.

Ngay lúc ấy\dots

*BÍP* BÍP*

Đồng hồ của tôi nhấp nháy. Tôi bước sang một bước.

*BỐP!*

Cú đấm của hắn nện thẳng vào mặt tên trùm.

``\textbf{Mày bị ngu à, thằng chó này!!}'' Tên trùm vừa ôm mặt vừa đấm ngược vào đầu hắn.

Thằng vừa ăn đòn ôm đầu: ``\textbf{Tao đâu có cố ý! Nó né nhanh quá!}''

Tôi đứng ngoài, khoanh tay, nhìn cảnh tượng hỗn loạn đó như thể đang xem hài kịch. Một tên còn lại run run chỉ tôi: ``\textbf{Mày\dots\ mày dùng ma thuật gì đó hả?}''

Tôi nhún vai: ``Ma thuật gì đâu, chỉ là phản xạ nhanh thôi mà.'' Cả bọn trố mắt, không hiểu nhưng cảm thấy sợ.

Một tên trong đó hét lớn, giơ nắm đấm lên cao: ``\textbf{All-in đi! Tổng tiến công!}''

Tên trùm quát lại: ``\textbf{Được! Đánh cho nó không kịp thở!}''

Rồi khi cả bọn bắt đầu gom lại chuẩn bị tấn công tổng lực, thì\dots

*BỤP!*

Một cú lên gối cực gọn đánh thẳng vào bắp đùi tên cầm đầu. Không mạnh, nhưng đủ khiến hắn tê rần cả chân.

``Bọn bây rảnh rỗi quá hả?'' tiếng Shiina-san vang lên, sắc lẹm như sao. Cô nàng bờm sư tử tung thêm hai cú đá rất nhanh, rất chính xác về phía hai thằng khác, khiến chúng bật lùi.

Ngay lúc đó, Nanase-san bước lên đứng chắn trước tôi, giọng đầy uy lực: ``Ở đây có camera đấy. Muốn tôi gọi thầy giám thị xuống không?''

Nhìn thấy hai người, cả đám lập tức đông cứng.

Fukami Shiina - chị đại thật sự của trường. Trận nào không đánh thắng thì cũng đánh hoà, không chột thì cũng què.

Furukawa Nanase - con gái nhà Furukawa, gia đình đã tài trợ và xây dựng phần lớn thị trấn này, trong đó có ngôi trường Aogami. Ông nội của cô là hiệu trưởng của trường, và bản thân gia đình cô cũng có số có má trong khu vực.

Chọn gây chuyện với tôi thì còn có thể, còn gây chuyện trước mặt hai người này thì\dots\ xin chia buồn. Họ gần như không có cửa thắng, không cẩn thận thì còn đi về chân lạnh toát.

Tên trùm lắp bắp: ``Thôi\dots\ anh em, rút! R-rút! Nhưng mày nhớ mặt tao đó, thằng mặt chó Hikawa kia! Lần tới tao sẽ đập mày ra bã!''

Tôi khoát tay: ``Ừ, nhớ thì nhớ. Nhưng đừng xuất hiện trước mặt tôi lúc chưa tỉnh táo đâu đấy.''

Chúng nguyền rủa thêm vài câu rồi bỏ chạy. Tôi chẳng buồn để tâm.

Ngay khi đám du côn rời đi, Yae-san và Hotaru-san chạy đến. Nhìn thấy tôi trông hơi xộc xệch so với ban nãy, cô nàng da ngăm trố mắt ngạc nhiên: ``Kuon-kun!? Chuyện gì vừa--''

Shiina-san bình thản đáp, vừa bẻ khớp tay răng rắc: ``À, có đám rảnh quá muốn nộp mạng cho tôi thôi mà.''

Cô bạn Cơ bắp tái mặt, còn cô bạn Da ngăm im luôn. Có vẻ như họ cũng biết rõ cô nàng bờm sư tử kia ghê gớm như thế nào.

\ct

Cả bốn người hộ tống tôi về lớp 1-C như thể tôi vừa thoát khỏi vụ ám sát. Trên đường đi, cô nàng bờm sư tử hỏi: ``Này Kuon-san. Sao cậu không né đám đó từ đầu đi?''

Tôi thản nhiên: ``Tôi biết mà. Kiểu gì mấy ông đó cũng giở trò. Tính cả rồi.''

Cô nàng tóc xanh phì cười: ``Thế hai chị em nhà cậu đâu?''

Tôi đáp luôn: ``Đang ngủ phình bụng ở nhà. Giờ này chắc mới bật dậy rồi đang hấp ta hấp tấp đây.''

Yae-san lập tức mặt méo xệch, thất vọng: ``Ngủ quên ư!? Mà Kuon-san bỏ họ ở nhà như vậy cũng quá đáng đó\dots''

Nanase-san kiêu kỳ đáp lại, như thể đó là điều hiển nhiên: ``Phải có lần nhớ đời thì mới chịu sửa.''

Hotaru-san thì gật gù đồng tình: ``Đúng là cần thiết thật. Cho chừa cái tội ngủ dậy muộn.''

``Tôi không muốn nghe vậy từ một người cú đêm như cậu đâu, Hotaru-san,'' tôi buột miệng trêu, khiến cô bạn nghệt ra mặt.

Shiina-san bụm miệng cười khoái trá: ``Heh. Hình dung cảnh đó xem. Có vẻ vui đấy.''

Khi về gần tới lớp 1-C, tôi tự thấy khóe miệng mình nhếch nhẹ.

Tôi có thể tưởng tượng được cảnh hai chị em nhà kia lúc này đang chạy khắp phòng, tóc rối tung, tìm đồng phục trong tuyệt vọng.

Cũng đáng mà. Gọi ba lần mà chẳng thèm dậy, bị đi muộn thì ráng phải chịu.

\pov{Hikawa Kaigou}

``\dots Ư\dots''

Tôi hé mắt. Ánh sáng từ cửa sổ chiếu thẳng vào mặt khiến tôi phải nheo lại. Đầu óc thì mụ mị, thân người thì nặng như chì. Tôi với tay ra theo thói quen, và tôi cảm thấy\dots

\dots Kuon-kun không còn ở đây nữa.

``Ơ\dots?''

Tôi bật dậy luôn, mắt mở to. Chẳng có cái lưng nào, chẳng có cái bóng cao to đang càu nhàu mỗi sáng nào bên cạnh.

``\textit{Kuon-kum\dots\ cậu ta đi đâu rồi?}''

Tôi quơ lấy cái điện thoại đang sạc trên đầu giường và bật lên. Đồng hồ hiện rõ chỉ 6:30.

\dots

\dots

``\textbf{\vie{\dots Chết thật rồi!!}}''

Tôi ngủ quá giờ rồi! Khỉ thật, tại sao cậu ta lại không gọi chúng tôi dậy chứ!

Chưa kịp suy nghĩ thêm điều gì, tôi đã lập tức lay dậy cô em gái vẫn còn đang quấn chăn quanh mình. ``\textbf{Suzu! Suzu! Dậy! Muộn giờ rồi!}''

``Ưm\dots'' Suzu khẽ cựa quậy, nói mớ: ``\vie{Onee-sama\dots\ Cho em ngủ thêm năm phút nữa\dots}''

``\textbf{\vie{Không có năm phút nào nữa! Sắp đến giờ vào học rồi!!}}''

Nghe tới đây, chỏm tóc ngố của Suzu dựng đứng lên. Con bé giật bắn dậy, tóc xù như ổ quạ, mặt ngơ ngác như đang ở hành tinh khác.

Hai chị em lao vào guồng xoay quen thuộc, nhưng hôm nay thì chẳng khác nào như bị tẩu hỏa nhập ma.

Suzuki ôm đồ chạy vào nhà vệ sinh trước, đóng cửa *CẠCH* một tiếng.

Một phút\dots

Hai phút\dots

\dots

``\textbf{\vie{Suzu! Em làm cái gì trong đó vậy!?}}'' tôi bắt đầu mất bình tĩnh, hét lớn từ bên ngoài.

``\vie{B-Bình tĩnh đi nào, Onee-sama! Em còn phải đánh răng nữa chứ\dots}'' giọng nói khẩn trương từ phía bên cánh cửa cất lên, rõ ràng cũng vội vàng không kém.

``\textbf{Đánh răng kiểu gì mà lâu thế!? Mở cửa ra để chị rửa mặt chứ!!}''

Bình thường, tôi sẽ không dễ dàng mất bình tĩnh như vậy, thậm chí còn khá kiên nhẫn. Nhưng việc chúng tôi đang phải chạy đua với thời gian, rõ ràng không thể tránh khỏi việc này.

Tôi đập cửa thình thình như sắp phá luôn bản lề. Trong khi đó, Suzu từ bên trong la oai oái: ``\textbf{\vie{Em xong nhanh ngay đây!}}''

Cuối cùng khi nó mở cửa, cả hai đổi chỗ nhanh tới mức tưởng đang ở một cuộc đua tiếp sức. Bàn chải, nước bắn tung tóe, tiếng nước chạy rào rào, tiếng hai chị em quát với nhau liên tục, đúng nghĩa hấp tấp khi mà do chúng tôi dậy muộn.

Tôi vừa lau mặt vừa chạy về phía tủ quần áo, giật phăng cái khăn xuống sàn: ``\vie{Suzu, áo sơ mi của em đâu rồi? Đừng nói là vứt lẫn vào chồng đồ bẩn nhé!}''

``\vie{Em\dots\ em không thấy!}'' giọn của em tôi loạng choạng. ``\vie{Onee-chan, lấy giùm em cái váy trong ngăn kéo thứ hai!}''

Tôi lục tung ngăn kéo, quăng cái váy về phía cửa: ``\vie{Bắt lấy! Mau mặc rồi xuống bếp!}''

``\vie{Chờ chút, khoá kéo kẹt rồi!}''

``\textbf{\vie{Giật mạnh lên!}}'' Tôi hét lại, tay không ngừng kéo phéc-mơ-tuya cặp sách. ``\vie{Này, sách Toán của chị đâu? Suzu, em có thấy quyển bài tập hôm qua không?}''

Suzu lao ra, váy còn nửa chừng khuy, tóc tai rối như tổ quạ: ``\vie{Em để trên bàn học\dots\ khoan, em lấy nhầm vào cặp của em mất rồi!}''

``\textbf{Trả đây!}'' Tôi giật lại quyển sách, nhét vội vào cặp. ``\vie{Khóa cửa phòng chưa?}''

``\vie{Em quên!}''

``\textbf{Chạy lên khóa!}'' Tôi đẩy em nó quay ngược, tay còn kịp nhét thêm hai chiếc bút vào túi váy. ``\vie{Khóa xong nhảy cầu thang ba bậc một!}''

Tiếng Suzu lục đục trên lầu vang xuống: ``\vie{Xong rồi!}''

Tôi vừa kéo khóa cặp vừa chạy ra ngoài, suýt va vào thành cầu thang: ``\vie{Xuống ngay!}''

Tôi vừa thở vừa rít lên khi nhảy ba bậc cầu thang một lúc. ``Kuon-kun cái tên khốn kiếp! Sao cậu ta lại bỏ hai chị em mình lại chứ!?''

Suzu vừa chạy vừa nấc: ``\vie{Onee-sama, em lo quá! Mình có kịp ăn sáng không đây?}''

``\textbf{Chị không biết!}'' tôi nghiến răng, mắt nhìn về phía bếp như thể sợ nó sẽ biến mất. ``Cậu ta đi một mình thì được, còn mình thì đói meo! Không lẽ phải nhịn đến trưa!?''

Tôi loạng choạng đáp xuống nền gạch, tim đập thình thịch vì vừa chạy vừa tưởng tượng cảnh bụng réo giữa giờ học. ``Kuon-kun đúng là đồ độc ác! Biết chúng ta dậy muộn mà vẫn không gọi, lại còn dám bỏ đi trước!''

Suzu rưng rưng: ``\vie{Em đói bụng rồi\dots}''

``\textbf{Khỉ thật!}'' tôi lại chửi thêm một lần nữa, khoát tay như muốn xua tan cái đói đang bò lên cổ họng. ``Nếu hôm nay mình ngất xỉu vì đói thì tôi sẽ\dots\ sẽ\dots\ giận cậu ta cả tuần!''

Chạy xuống dưới bếp, hai chúng tôi khựng lại khi thấy hai phần bữa sáng được gói gọn gàng, có thể vừa cầm ăn vừa chạy. Đi cùng với đó là hai hộp cơm trưa đặt chồng lên nhau, xếp ngay ngắn ở trên bàn, cùng với hai tờ giấy nhỏ được gắn trên chúng. Nét chữ này có lẽ là từ Kuon-kun, tôi có thể chắc chắn.

Hộp cơm của tôi có ghi:

\txtbx{``Có muộn đến mấy cũng phải ăn, kẻo học giữa chừng đói lả. Nhớ khóa cửa hộ tôi.''}

Trên hộp của Suzu, cậu ta viết:

\txtbx{``Nếu em để đói rồi xỉu thì đừng đổ thừa. Lần sau cố mà dậy sớm.''}

Lật mặt tờ giấy của tôi ra đằng sau, cậu ta còn viết thêm:

\txtbx{``P.S\: Tôi đoán chắc chắn hai người sẽ dậy muộn, nên làm loại cầm tay.''}

Tôi siết tờ giấy, mặt nóng lên vì tức xen lẫn xấu hổ. ``\vie{Cậu ta biết trước luôn\dots}''

Suzu vừa nhét hộp cơm vào cặp vừa lẩm bẩm: ``\vie{Nhưng\dots\ cũng biết chúng ta dậy muộn mà\dots}''

Hai chị em nhìn bữa sáng, nhìn hai mảnh giấy, rồi nhìn nhau.

Và rồi\dots

``\textbf{Onii-sama/Kuon-kun xấu tính!!!}''

Tất nhiên đó chỉ là lời ``nguyền rủa'' kiểu trách yêu, vì lỗi này rõ ràng là do chúng tôi.

\dots Ít nhất thì, cậu ta còn chuẩn bị trước cho chúng tôi. Ơn trời.

\ct

Cả hai lao ra khỏi nhà như đang bị truy đuổi, tay nắm chặt bữa sáng còn ấm mùi trứng cuộn và xúc xích, mùi mè rang vương trên ngón tay. Hai hộp cơm trong cặp thỉnh thoảng va vào, kêu *LỘP CỘP* như tiếng mõ báo hiệu thảm họa.

``\vie{Aaa! Chân em đau! Onee-sama, chậm lại đi!!}''

``\vie{Không chậm được!! Còn năm phút nữa là chuông reo rồi!!}''

``\vie{Em ghét cái đồng hồ báo thức\dots\ Em ghét luôn Onii-sama nữa\dots\ Cũng tại tối qua em ngủ quá ngon rồi\dots}''

Hai chị em vừa chạy vừa than, vừa tự đổ lỗi lung tung để giải tỏa sự bối rối vì dậy trễ.

Gió sáng hun hút thổi ngược, xoáy tung tóc tôi thành hai cánh hoa phai màu hồng, mồ hôi lấp lánh trên thái dương Suzu như những hạt sương nhỏ xíu vừa kịp nở. Cổ họng tôi cay xè mùi xúc xích nguội, nhưng không ai dám dừng lại để nhai; chỉ kịp ngậm một miếng cơm mè, nhai tóp tép rồi nuốt ừng ực, sợ hạt cơm văng ra đằng sau như tàn pháo.

Bóng chúng tôi in trên nền đường dài và loang, hai cái cặp sách nảy lên nảy xuống như những con cá mắc cạn, dây kéo va vào nhau leng keng. Một con chó hoang giật mình nhảy vội vào lề, nhìn theo bóng chúng tôitrẻ biến mất trong sương bụi cuối phố, chỉ còn tiếng hộp cơm loảng xoảng và tiếng thở gấp như bản nhạc jazz bị kéo nhanh gấp đôi.

\ct

Vừa vào cổng trường, tôi và Suzu gần như phóng nước rút đến hành lang lớp. Giày đập xuống nền gạch vang lên lách cách như mưa rào, cái cặp sách cứ nhảy tung trên lưng, suýt nữa tuột dây. Tôi thở gấp, mồ hôi đã bắt đầu thấm sau cổ áo, nhưng chẳng dám giảm tốc, muộn thêm một giây nữa thì coi như trễ học.

Băng qua khúc quanh, tôi giật mình thắng lại kịp thời. Phía trước, đúng một đám con trai đang lảo đảo bước ngược chiều. Thằng cầm đầu còn ôm hông, mặt méo xệch, trông như thể vừa trải qua một trận đánh. Chúng nghe tiếng bước chân gấp, đồng loạt ngẩng lên.

Một tên đằng sau trố mắt, chỉ thẳng vào tôi: ``Nhìn kìa, đại ca! Lại là cái tên Hikawa Kuon đó kìa!''

``Nhưng mà\dots\ sao bây giờ nó lại trông khác vậy? Lại còn dẫn nhỏ nào đó nữa kia?''

Tôi khựng lại nửa bước. Hóa ra chúng nhầm tôi với Kuon-kun vì nét mặt và dáng đi giống cậu ta. Mà chuyện gì đã xảy ra giữa chúng với cậu ta vậy?

Trong đầu tôi lập tức tính toán: nếu dừng lại giải thích thì mất thời gian, mà chúng đang bực tức, khéo còn bị vạ lây. Thôi thì, chốc nữa hỏi cậu ta một tiếng về sau vậy.

Tôi không đợi chúng kịp nói thêm câu nào.

``Tránh. Ra.''

Một cú đấm thẳng. Gọn. Nhanh. Mạnh.

Đến mức tên đứng đầu bay nguyên cụm, kéo theo hai thằng sau suýt ngã chồng lên nhau. Cả bầy đập vào tường tạo tiếng *RẦM!!* vang cả tầng, khiến mấy học sinh gần đó giật mình quay lại.

Tên trùm ngã gục, bất tỉnh nhân sự. Đám tay sai hoảng loạn, có đứa run run lay hắn, có đứa lùi dần về phía sau, mắt vẫn không dám rời khỏi tôi.

Tôi hét lớn, giọng vang vọng hành lang: ``\textbf{Chúng tôi đang vội! Cút ra!!}''

Suzu chạy sát phía sau, thở không ra hơi, cúi đầu lia lịa: ``Xin lỗi! Xin lỗi các anh nhiều lắm!! Bọn em thật sự đang gấp!''

Chúng tôi lao tiếp, bỏ lại phía sau đám côn đồ ngơ ngác. Tim tôi đập thình thịch, nửa vì sợ muộn, nửa vì hồi hộp, đây là lần đầu kể từ ở The Void, mình ra tay như vậy. Nhưng giờ không phải lúc nghĩ ngợi; còn hai tầng nữa mới đến lớp.

Vừa rẽ vào cầu thang bên phòng, tôi thoáng thấy một nhóm sáu học sinh đang tụ tập trước phòng giáo viên. Trông họ lạ mặt, cặp sách còn dán tem mới tinh, đồng phục gấp nếp thẳng thớm. Có lẽ là các bạn chuyển trường? Tôi không kịp dừng lại để xác nhận.

``Nhanh lên, Suzu!'' tôi quay sang em gái, kéo tay cô bé chạy vượt lên ``Còn hai phút nữa thôi!''

Chúng tôi lao lên lầu, tiếng giày loẹt quẹt dồn dập như trống trận. Trong đầu tôi chỉ còn một mục tiêu duy nhất: kịp giờ vào lớp, trước khi chuông reo.

\ct

*REEEEEEENG*

Tiếng chuông trường vang lên đúng lúc hai chị em chạm tới cửa lớp.

Tôi mở xoạch cánh cửa ra, lao tới ngay giữa bục giảng, hướng về chiếc bàn giáo viên, chẳng hề để ý xung quanh.

Cả hai đồng thanh hét lớn: ``\textbf{Xin lỗi thầy, bọn em đến muộn!!!}''

Sau đó\dots

\dots

Im lặng.

Ba giây sau, cả lớp phá lên cười. Trong lớp chẳng có thầy cô nào cả. Đây vẫn là giờ truy bài.

Suzu đứng đơ, mặt đỏ như cà chua. Tôi thì vừa xấu hổ, vừa tức, vừa muốn độn thổ vì bao nhiêu tiếng la hét chạy dọc hành lang cuối cùng thành trò cười.

Tôi quay sang nhìn Kuon-kun, cái người đang ngồi mở sách, mặt tỉnh bơ như thể đã biết mọi thứ.

``Onii-sama\dots'' Suzu phồng má, tiến về bàn của cậu ta: ``Anh thật xấu tính\dots''

Kuon-kun ngẩng lên, thở dài: ``Anh gọi hai người tận ba lần. Hai người không thèm dậy. Thế là thôi, đành phải đi trước.''

Cả lớp lại rộ lên một tràng ``Oaaaaaa\~{}'', Hotaru-san thậm chí còn buột miệng trêu: ``Là lỗi hai cô thật mà.''

Saya-san thì nhìn chúng tôi, không nhịn được mà cười: ``Nhưng mà cách hai cậu la hét ở ngoài kia nghe dễ thương ghê\~{} Nhất là bản mặt ngố tàu khi lao vào lớp đó\~{}''

Tôi vừa định mở miệng phản bác cậu ta thì Nanase-san đã lên tiếng, giọng nhẹ nhưng đủ để cả lớp nghe rõ: ``Kuon-san, dù sao cũng nên chờ hai người họ thêm chút xíu. Đồng hồ sinh học cần thời gian để điều chỉnh mà.''

Miku-san gật gù, tay gõ nhẹ bàn: ``Đúng vậy. Hơn nữa, để hai cô bé chạy như vậy trên hành lang rất nguy hiểm. Lần sau, nếu không dậy được, cứ gọi tôi, tôi sẽ mang nước đá đến tận giường.''

Tôi đỏ bừng, vừa định cãi lại thì Kuon-kun lại thở dài, đóng sách lại: ``Được rồi, được rồi. Lần sau tôi sẽ để lại đồng hồ báo thức thật lớn cho hai người. Nhưng mà nếu còn ngủ nướng, anh sẽ chụp ảnh lại rồi đăng lên group lớp đấy.''

``\textbf{Đừng mà!!}'' cả hai chị em hét lên, khiến cả lớp lại cười ầm.

Tôi lúc này chỉ muốn tìm cái hố nào chui xuống thật sự\dots\ Thật không thể tin được cái buổi sáng éo le này lại đến mức này\dots

\pov{-}

Khung cảnh lớp 1-C sau tràng cười dần trở lại trạng thái bình thường. Tiếng sột soạt sách vở, tiếng bút lách cách, đâu đó vài tiếng thì thầm trao đổi bài tập. Cả lớp như vừa trở lại guồng quay học tập quen thuộc.

Ở khu vực gần bàn giáo viên, Kuon ngồi giữa, hai chị em nhà Hikawa nghiêng về phía cậu. Cậu Tổng lật sách, chỉ tay nhẹ vào bài, giọng trầm ấm: ``Công thức này dùng cho bài 3, đừng nhầm với bài 2.'' Cô chị lớn gật nhanh, mắt vẫn còn vương ngái, còn cô em gái cặm cụi ghi chú, má ửng hồng vì cố gắng tập trung.

Shino ngồi một mình bàn bên, lưng thẳng tắp, tay trái đỡ cằm, tay phải lật trang sách nhịp nhàng như máy. Ánh đèn huỳnh quang phản chiếu trên tóc đen dài, môi cô mấp máy đọc thầm, hoàn toàn chìm trong thế giới riêng.

Dãy giữa, Miku và Kaoru dựa đầu vào nhau. Cô nàng tóc ngắn thì thào: ``Đoạn này tôi dịch được rồi, cho tôi xem đi.'' Cô nàng bờm sói gõ bút lên trán cô: ``Xem để rồi để chép à? Tự làm đi, mai thầy thu đấy.'' Dứt lời vẫn đẩy tờ giấy có phác thảo sang, mắt trao nhau nụ cười khẽ.

Shiina ngồi bàn cuối, chân gác lên ngăn bàn, tay xoay bút như chong chóng, mắt trừng lại quan sát xung quanh. Honoka bên cạnh cầm bánh quy Wata chia, nhỏ giọng: ``Shiina-san, cậu có muốn ăn kẹo dẻo không?'' Cô bạn bờm sư tử nhếch môi: ``Không rảnh. Cậu câm miệng lại đi.'' Nhưng vẫn giật lấy một viên kẹo, nhai rào rào.

Kanade và Kana ở hàng cửa sổ. Cô nàng thiên sứ vẽ hình khối trong vở, thỉnh thoảng chỉ cho cô bạn thân cách giải bài tập Hình. Cô bạn ác ma gật gù, mắt lim dim vì ánh nắng sớm, tay vẫn chép theo từng nét bút của chị, mái tóc nâu đỏ đung đưa theo gió lùa qua khe cửa.

Từ phía cửa lớp, tiếng giày cao gót nhịp nhẹ xuống sàn gốm vang lên những tiếng *CẠCH CẠCH*, khiến cả lớp 1-C - vốn vừa ổn định sau màn chạy nước rút của hai chị em nhà Hikawa - đồng loạt ngẩng đầu nhìn.

Một người phụ nữ bước vào.

Mái tóc vàng-đỏ mềm mại đổ xuống vai, gợn sóng nhẹ theo mỗi bước chân. Ánh mắt dịu dàng, nhưng sâu trong đó lại ẩn một nét cứng rắn khó diễn tả. Thần thái của cô khiến toàn lớp lập tức ngay ngắn hơn bình thường, kiểu khí chất ``mềm mỏng như người mẹ, nhưng có gì đó nghiêm khắc'' mà chỉ vài người mới toát ra được.

Trong mắt Kuon, cô khiến cậu lập tức liên tưởng đến Yuzuki Choco, một nữ y tá ở Ác giới, nhưng là một phiên bản hoàn toàn con người, không mang đặc điểm gì của quỷ, không chút gì siêu nhiên. Chỉ đơn thuần là một phụ nữ đẹp, gợi cảm theo cách rất nền nã và chỉn chu, với áo xếp nếp màu xanh lá, đi với một cái váy bút chì màu đen (về cơ bản, không khác mấy so với cô nàng Y tá ở thế giới gốc).

Kaigou, Suzuki và một vài học sinh khác khẽ liếc nhìn nhau. Đúng như giáo viên tạm quyền hôm qua thông báo, giáo viên chủ nhiệm thật sự sẽ đến vào hôm nay. Và họ gần như chắc rằng đây chính là người đó.

Người phụ nữ đứng trước bục giảng, đặt túi xuống bàn, rồi mỉm cười với lớp.

``Tôi là Tsukishita Yuko.'' Giọng cô dịu mà rõ ràng, từng chữ như được đặt đúng vị trí. ``Tôi sẽ là chủ nhiệm của lớp 1-C kiêm giáo viên dạy Sinh học trong năm nay.''

\chrint

Tsukishita Yuko (\ruby{月}{つき}\ruby{下}{した}\ruby{癒}{ゆ}\ruby{子}{こ}, \ruby{Xư-ki}{Nguyệt}-\ruby{si-ta}{Hạ}\ \ruby{Giu}{Dũ}-\ruby{cô}{Tử}) đảm nhiệm vai trò giáo viên chủ nhiệm lớp 1-C và giảng dạy bộ môn Sinh học của trường Trung học Aogami.

Yuko dành cảm tình đặc biệt cho những học sinh siêng năng, thường xuyên biên soạn thêm tài liệu ngoài giáo trình hoặc dành thời gian trợ giúp cá nhân cho bất kỳ ai thể hiện quyết tâm học tập. Các bài giảng của cô trở nên sinh động nhờ những ví dụ thực tế từ thiên nhiên: từ cách loài bọ hung cuộn phân đến chu kỳ phá hủy và tái sinh của rừng mưa, giúp học sinh dễ hình dung và ghi nhớ lâu. Tuy nhiên, với những kẻ lười biếng, cô trở nên nghiêm khắc: những bài kiểm tra bất ngờ có thể xuất hiện bất cứ lúc nào nhằm ``giữ nhịp tim cả lớp tỉnh táo'' (theo lời phân trần của một học sinh đã từng được dạy bởi chính giáo viên này.).

Thầm lặng, Yuko mang trong lòng nỗi lo về lịch sử kỳ lạ của trường Aogami; mối lo âu ấy thỉnh thoảng lẫn vào bài giảng qua những "lời cảnh báo giả định" như câu chuyện loài cá biến mất khỏi hồ nước bị bỏ hoang. Trong mắt Hikawa Kuon, người từng quen biết Sakura—một nhà môi trường học ở vòng lặp trước—có lẽ giáo viên này sẽ cực kỳ "hợp cạ" với Sakura, cả về khía cạnh khoa học lẫn tinh thần bảo vệ thiên nhiên.

\chrint

Cả lớp im lặng, chờ những lời tiếp theo.

Yuko đặt hai tay lên bàn, hơi nghiêng đầu, nụ cười vẫn dịu nhưng ánh mắt lại sắc hơn một chút.

``Trước khi bắt đầu, tôi sẽ nói rõ hai điều về cách tôi dạy học.''

Cả lớp lập tức thẳng lưng.

``Một: \emph{tôi rất ghét những học sinh cẩu thả và lười biếng.} Tôi sẽ không trừng phạt ngay đâu, nhưng tôi nhớ rất lâu đấy.'' Một số học sinh nuốt nước bọt. Kaigou thì khẽ rùng mình, còn Suzuki thì ngồi thẳng tắp, hai tay đặt đúng tư thế như học sinh gương mẫu.

``Hai: \emph{tôi có thói quen kiểm tra đột xuất.} Không báo trước, không nhắc nhở. Các em nên quen với điều đó từ bây giờ.'' Cô nói câu đó bằng giọng nhẹ như gió, nhưng lại có chút tinh nghịch như thể rất mong chờ được ``bắt bài'' ai đó.

``Và ba: \emph{tôi biết hết tất cả những gì mà các em làm ở bên dưới.} Đừng hòng qua mặt được tôi.'' Trong khi cả lớp căng thẳng ra mặt, Kuon chỉ nhẹ nhàng thở ra. Đối với cậu, kiểu giáo viên này quá quen thuộc.

Từ cấp hai, cấp ba cho tới khi lên đại học, cậu từng gặp không ít những người như thế: nghiêm khắc, có nguyên tắc, đôi khi hơi ``đáng sợ,'' nhưng lại rất dễ đối phó. Chỉ cần trở thành một học sinh mẫu mực - điều mà Kuon vốn đã làm một cách tự nhiên - là đủ. Không cần diễn.

``\textit{Cái kiểu này\dots\ Suzuki chắc cũng ổn thôi.  Con bé trông cũng khá gương mẫu lễ phép. Nhưng Kaigou-chan thì\dots}'' Kuon nhìn sang cô chị đang len lén cúi đầu. ``\textit{\dots xem chừng không chắc lắm. Mình còn chẳng biết cô ấy hồi ở thế giới cũ học hành kiểu gì nữa.}''

Những bạn trong lớp học bắt đầu xì xào.

Uzuki đưa tay che miệng, nói khẽ với người bên cạnh: ``\hl{Nghe như kiểu hiệu trưởng của học viện quân sự vậy.}''

Kaoru khoanh tay, nhướng mày: ``\hl{Không tệ. Ít nhất còn nghiêm túc hơn mấy ông thầy già hay ngủ gật. Suzuki-san, cậu nghĩ sao?}''

Suzuki thì gật đầu liên tục, đôi mắt sáng như sao: ``\hl{Mình thích giáo viên nghiêm một chút! Như vậy em mới chăm chỉ hơn được!}''

Kaigou quay qua nhìn em gái mình với ánh mắt không thể tin nổi, tựa như muốn hỏi ``em thật sự sống ở thế giới nào vậy\dots?'', còn một số bạn khác thì xen lẫn cả hào hứng lẫn lo lắng.

Yuko mỉm cười trước phản ứng của lớp, như thể hài lòng với sự hỗn hợp giữa căng thẳng và mong đợi ấy. ``Được rồi,'' cô vỗ tay một tiếng. ``Ta vân còn chút thời gian trước khi vào giờ. Vậy thì tôi sẽ giới thiệu những bạn học mới sẽ học cùng lớp chúng ta.''

Cả lớp hướng mắt ra cửa. ``Mời sáu em vào.''

Từ ngoài cửa, sáu bóng người bước vào sau lời mời của cô giáo. Ngay khoảnh khắc họ xuất hiện, cả lớp đồng loạt ồ nhẹ, một phần có lẽ bởi vì bất ngờ.

Kaigou chớp mắt, bỗng thấy tim đập loạn nhịp. ``\textit{\dots Mấy người này\dots\ chẳng phải là\dots}'' Cô lập tức nhận ra họ chính là nhóm đã đứng chờ trước phòng giáo viên lúc nãy, khi cô vội vã chạy ngang qua, chỉ kịp lướt mắt thấy bóng dáng họ dưới ánh đèn huỳnh quang

Riêng Kuon thì hơi tê sống lưng. ``\textit{\dots Không thể nào. Sáu người này\dots\ trông giống đến khó tin.}'' Dĩ nhiên, không có sừng, không có đuôi, không có đặc điểm phi nhân như ở thế giới gốc. Nhưng gương mặt, phong thái, thần thái toát ra\dots\ tất cả đều quá giống.

Yuko gật đầu, cho người đầu tiên bước ra phía trước.

Cô gái đầu tiên bước lên theo lệnh. Cô có mái tóc đen ngắn kiểu bob, đeo một chiếc nơ ở bên trái, dáng nhỏ nhắn nhưng ánh mắt tím mạnh mẽ. Khi cúi đầu chào, cô cúi hơi\dots\ quá sâu, gần như gập cả người.

``Chào\dots\ chào mọi người! Mình là Utachi Azuki. Mình thích những chuyến mạo hiểm và\dots\ ừm\dots\ mình cũng hát khá tốt. Mong được giúp đỡ!''

\chrint

Utachi Azuki (\ruby{歌}{うた}\ruby{地}{ち}\ruby{亞}{あ}\ruby{槻}{ずき}, \ruby{U-ta}{Ca}-\ruby{chi}{Địa}\ \ruby{A}{Á}-\ruby{dư-ki}{Quy}) là một cô gái mang trong mình tinh thần mạo hiểm bẩm sinh. Từ nhỏ, cô đã bị cuốn hút bởi những điều mới lạ và không ngần ngại thử thách bản thân trong những tình huống khó lường. Không chỉ dừng lại ở việc khám phá, Azuki còn sở hữu một giọng hát đầy nội lực, có khả năng chạm tới trái tim người nghe chỉ qua vài câu ca đầu tiên. Chính sự kết hợp giữa sự liều lĩnh và tài năng âm nhạc đã khiến cô trở thành một nhân vật đặc biệt trong mắt bạn bè.

Không ai có thể phủ nhận niềm vui sống mãnh liệt của Azuki. Cô luôn tìm kiếm những khoảnh khắc đầy phấn khích và thường xuyên kéo những người bạn thân nhất vào những cuộc phiêu lưu bất ngờ. Từ việc thám hiểm các tòa nhà bỏ hoang nằm rải rác quanh thị trấn Aogami, nơi mà người dân địa phương thường tránh xa vì những lời đồn thổi, cho tới việc tổ chức những buổi biểu diễn ngẫu hứng ngay tại trung tâm thị trấn, nơi cô và nhóm bạn vừa hát vừa múa, thu hút sự chú ý của mọi người. Với Azuki, không có khái niệm ``không thể'' chỉ có ``chưa thử''.

Trong câu lạc bộ âm nhạc của trường Aogami, Azuki không chỉ là thành viên, cô còn là linh hồn. Với vai trò ca sĩ chính, cô đứng trên sân khấu như một ngọn lửa sống, thiêu đốt mọi sự do dự và truyền cảm hứng cho cả nhóm bằng năng lượng không sợ hãi. Những bản ballad sâu lắng hay những ca khúc rock mãnh liệt đều được cô thể hiện một cách chân thật và đầy cảm xúc, khiến người nghe không thể rời mắt. Dưới ánh đèn sân khấu, Utachi Azuki không chỉ hát, cô kể chuyện bằng âm thanh, và mỗi câu chuyện đều bắt đầu bằng một bước đi đầy mạo hiểm.

\chrint

Cả lớp bật cười nhẹ vì sự vụng về nhưng đầy nhiệt huyết ấy.

Kuon nhìn cô, lòng khựng lại một nhịp. Tên của cô nàng đó có cách đọc khá giống với AZKi, cô nàng diva của thế hệ thứ Không, cũng là thành viên mà cậu đã tìm hiểu qua bản sơ yếu, nhưng chưa gặp mặt trực tiếp lần nào.

``\textit{Azuki-san\dots? Nếu mình nhớ không nhầm thì ba tháng nữa em ấy mới lên văn phòng từ sau vụ giải tán của INNK nhỉ\dots?}''

Một cảm giác vừa thân thuộc vừa chua xót kéo ngang ngực cậu.

``Được rồi, người tiếp theo lên đi,'' Yuko nói.

Một cô gái tóc vàng với vài sợi dây đỏ cuốn quanh, vẻ tự tin nhưng hơi gắt. Cô bước ra, khoanh tay, ``hừm'' một tiếng như thể đang kiêu ngạo những người ngồi dưới. Cô đứng trước bục giảng, tay khoanh lạnh lùng, đôi mắt màu hổ phách nhìn xuống cả lớp như thể đang đánh giá từng kẻ lười biếng ấy.

``Tôi là Akaba Koishin. Các cậu đừng có mà hiểu nhầm; tôi không thân thiện gì đâu.''

Nói xong, mắt cô đảo nhanh sang lớp xem có ai nhìn không, mặt hơi đỏ.

\chrint

Akaba Koishin (\ruby{赤}{あか}\ruby{羽}{ば} \ruby{恋}{こい}\ruby{心}{しん}, \ruby{A-ca}{Thích}-\ruby{ba}{Vũ}\ \ruby{Cô-i}{Luyến}-\ruby{sin}{Tâm}), cái tên nghe nhẹ nhàng, nhưng ẩn sau đó là một cô gái luôn mang theo bức tường lửa tsundere cao ngất. Cái kiểu câu thường thấy của cô nàng này là ``Đừng có hiểu lầm, tôi chỉ\dots\ hứng thú một chút thôi!'' vang lên trước khi cô vội vã quay đi, tai ửng hồng.

Trong lớp, Koishin tạo dựng hình ảnh một ``viên đạn bạc'' lạnh lùng: chân gác bàn, sách mở ra nhưng mắt lướt qua mọi người bằng ánh nhìn ``không đủ trình để nói chuyện''. Cô than thở về ``phiền phức tập thể'', thề sẽ ``không bao giờ tham gia kịch nghệ đường phố hay hoạt động nhóm'', thế nhưng cô lén lút giúp họ, rồi vào phòng chỉ nói gọn: ``Nhỡ tay làm vỡ, đừng tưởng tôi cố ý giúp.''

Văn chương là cửa sổ Koishin khép hờ. Cô viết những truyện ngắn đặt tên mơ màng đẫm nước mắt nhân vật, đầy ẩn dụ về kẻ cô đơn mang gai. Khi hội trường văn học tổng duyệt, cô đứng sau rèm, tay siết kịch kín kịch bản, mắt dán vào khán giả; lễ phát giải cô lẩn trốn ngoài hành lang, nhưng bản thảo đoạt giải lại được đặt cẩn thận trong hộp bìa cứng, kèm dòng chữ ``Tặng người đã từng đọc, đừng bao giờ nhận là tôi viết.''

Bên trong lớp vỏ tsundere ấy là một nỗi sợ yếu đuối ám ảnh từ ngày cô còn học tiểu học. Từ sự kiện đó, cô học cách mặc áo giáp gai: không tin tưởng, không thổ lộ, không khóc. Mỗi lần có ai đó vượt qua rào cản, cô lại giận dữ đẩy họ ra, rồi đêm xuống lại ngồi cạnh cửa sổ, ôm gối, mắt đỏ hoe thì thầm tự trấn an bản thân.

Dẫu vậy, ở tận cùng trái tim, Koishin vẫn nuôi một bông hoa nhỏ, một niềm tin rằng có người đủ kiên nhẫn đi qua bão gai, thấy được tia ngọt ngào cô giấu kỹ. Và nếu một ngày bông hoa ấy dám nở, cô sẽ không còn là tsundere đơn độc nữa, mà chỉ là Koishin, người viết truyện, kẻ mộng mơ, và cô gái biết yêu thương.

\chrint

Kuon bất giác nở nụ cười khó thấy. ``\textit{Haachama-chan phiên bản trường học\dots\ trông giống như hồi em ấy mới vào làm, khí chất tsun đúng là không chạy đi đâu được.}''

``Cảm ơn em, Akaba-san,'' Yuko gật đầu, dù rằng nét mặt cô hiện thoáng qua sự lo ngại qua cách mà cô nàng tóc vàng nói chuyện. ``Em tiếp theo lên đi.''

``V-Vâng ạ!''

Cô gái tóc hồng-xanh nhạt thoáng giật mình, bước ra như sắp vấp vào chính mình.

``Chào\dots\ chào mọi người, tớ là Mizuno Aku\dots\ Tớ hơi\dots\ hậu đậu, xin mọi người hãy giúp đỡ mình thật tốt!''

\chrint

Mizuno Aku (\ruby{水}{みず}\ruby{野}{の}\ruby{亜}{あ}\ruby{久}{く}, \ruby{Mi-dư}{Thủy}-\ruby{nô}{Dã}\ \ruby{A}{Á}-\ruby{cư}{Cửu}) là một cô gái hậu đậu nổi tiếng trong trường, người luôn trong tình trạng ``vừa chạy vừa ngã,'' ``động đâu hỏng đó''.

Cô thường xuyên bị cô bạn thân Shion dụ dỗ làm chân sai vặt: từ việc photo tài liệu câu lạc bộ, mang đồ uống cho mọi người, đến việc chạy lên xuống các tầng lầu để tìm một ``bản in bí ẩn'' mà Shion đột ngột thấy cần gấp.

Đương nhiên, trên đường đi cô không tránh khỏi việc tự vấp vào chân bàn, làm đổ cà-phê lên áo mình, hoặc đâm sầm vào cửa kính đã mở sẵn, những tai nạn nhỏ trở thành đặc sản mà mọi người đều biết đến.

Tuy nhiên, đằng sau vẻ ngoài lúng túng ấy lại là một game thủ hàng hiệu. Dưới biệt danh bí mật là ``\textsc{AkUtopia}'', cô thống trị các bảng xếp hạng trực tuyến, phản xạ cực nhanh và óc chiến thuật sắc bén đến mức khiến đối thủ phải nghi ngờ cô dùng ma thuật. % textsc: chữ hoa nhỏ

Chính kỹ năng ``nhìn thoáng qua là biết chỗ bí mật'' này đôi khi lại giúp Câu lạc bộ Huyền bí giải mã những câu đố khó nhằn theo cách không ai ngờ tới; ví dụ, cô chỉ cần nhìn lướt qua dãy ký tự lạ là lập tức nhận ra đó là\dots\ mật mã Konami. Vì thế, dù Aku luôn than thở ``Tại sao mình lại đồng ý làm việc này chứ\dots'', cả câu lạc bộ đều thầm thừa nhận:  ``Nếu thiếu cô ấy, chúng ta đã bị kẹt từ\dots\ vòng gửi xe rồi.''

\chrint

Cô vừa cúi đầu vừa liếc thấy Kuon, người mà cô vô tình đụng phải trên đường cô tới trường. Mặt đỏ bừng.

``Ơ-Ơ-Ơ\dots\ Là cậu ban sáng! Chuyện vừa rồi, tớ—tớ xin lỗi rất nhiều! Là lỗi của tớ, tớ bất cẩn, tớ--''

Yuko phải ho nhẹ nhắc cô tiếp tục tự giới thiệu. Cả lớp cười ồ.

Kuon che miệng cười khẽ. ``\textit{Y hệt Aku-tan, không khác một li nào.}''

``N-Nhân tiện thì,'' Aku hướng về giáo viên, rồi về phía cả lớp, giới thiệu cô bạn đứng cạnh mình, ``đây là Shitou Shion, bạn thân của mình.''

Cô gái tóc bạc nhỏ nhắn bước ra, đặt tay lên ngực như pháp sư tối thượng.

``Ta là Shitou Shion, phù thủy xứ Vương Quốc Bóng Tối Vĩnh Cửu, thủ lĩnh CLB Huyền Bí! Nhờ sức mạnh của pháp thuật bóng tối, ta đã thấu rõ mọi bí ẩn trong vũ trụ! Hỡi các linh hồn tân sinh, hãy cúi đầu trước ánh mắt của ta!''

Cô giơ hai tay lên cao, ngón trỏ và ngón giữa chụm lại tạo thành dấu hiệu kỳ bí, mắt nhắm nghiền như đang niệm chú.

``\dots và đừng quên nộp lệ phí tham gia CLB là 500 yên một tháng!''

Cả lớp bật cười ầm ĩ. Yuko đứng bên cạnh đưa tay lên trán như đang đau đầu, nhưng khóe miệng vẫn không kìm được một nụ cười bất lực.

\chrint

Shitou Shion (\ruby{紫}{し}\ruby{藤}{とう}\ruby{詩}{し}\ruby{音}{おん}, \ruby{Si}{Tử}-\ruby{tô-ư}{Đằng}\ \ruby{Si}{Thi}-\ruby{ôn(g)}{Âm}) không chỉ là một ``chuunibyou'' đơn thuần, mà còn là đấng ``Thủy Tổ Ma Ảnh'' tự phong, kẻ thề chỉ đôi mắt bất tử mới đủ chiêm ngưỡng Vương Quốc Bóng Tối Vĩnh Cửu nơi cô ngự trị - hoặc ít nhất là dưới góc nhìn của một người ảo tưởng sức mạnh như cô. Cô thường xuất hiện với bộ phụ kiện như áo choàng voan đen dài chấm gót được may bằng lụa tự nhuộm bằng vỏ dâu rừng Aogami, hay chiếc``Quyền Trượng Hắc Ám'', thứ thực chất được tạo từ một cây gậy bóng chày cũ được bọc băng keo đen và đính hạt cườm thủy tinh mua ở chợ đồ cũ.

Bên ngoài vỏ bọc mạnh mẽ ấy, Shion thực ra là thư ký duy nhất của CLB Huyền bí, người tự tay thiết kế ``Phòng Thí nghiệm Tường thuật'' trong căn phòng kho tầng hầm dưới thư viện cũ.  Tại đây, cô lưu giữ hàng trăm bản đồ tay vẽ các lối đi ngầm dưới Aogami, ghi chép những bản tường thuật bằng mực tím pha lẫn tro nhang, và treo một tấm bảng ``Không được quay lưng với bóng tối sau 23:00,'' quy định mà chỉ cô và Aku dám vi phạm.

Dù hay lôi Aku vào những ``nghi thức'' ngớ ngẩn như đứng giữa bãi cổ tự lúc nửa đêm để ``lắng nghe tiếng khóc của tháp chuông'', Shion lại là người đầu tiên phát hiện ra rằng những huyền thoại đô thị ở Aogami không chỉ là truyện ma đơn thuần một cách đầy trùng hợp tới khó tin.

\chrint

``\textit{\dots Nhìn như thể là Shion-chan khi bị mất đi sức mạnh của phù thủy nhỉ?}'' Kuon bất giác thì thầm, mắt mở to như vừa nhìn thấy quá khứ đuổi kịp mình.

Ngay khi Shion hoàn thành xong màn giới thiệu của mình, cô gái tóc trắng đỏ, với nụ cười rạng rỡ, cúi đầu chào rất lịch sự.

``Yo, Onidoro Ayame. Cũng là hội trưởng hội học sinh năm nay\~{}''

Và ngay khi ngẩng đầu lên, cô híp mắt, đưa ngón tay thon thả chọc nhẹ vào sườn bạn ngồi cạnh như một cú tấn công bất ngờ, khiến nạn nhân giật mình ``kya!'' và cả lớp bùng nổ tiếng cười sảng khoái.

Rồi, đôi mắt đỏ của cô nhanh như cắt lia thẳng về phía Kuon. Khóe môi cô nhếch lên thành một cung tàn nhẫn, nụ cười tinh nghịch đến mức đáng nghi.

``Ồ\~{} lớp này chỉ có mỗi một cậu con trai thôi à? Vui đấy ha\~{}''

Cô khẽ nghiêng đầu, giọng ngọt sớt như mật nhưng đầy gai: ``Cậu là Kuon-kun đúng chứ\~{}? Đừng có mà run nhé, tôi không cắn đâu. Trừ khi cậu muốn.''

Giọng chợ búa nhẹ, y như bản gốc, chỉ thêm một thìa đường và một nhúm ớt.

Kuon thở dài trong đầu. ``\textit{Ojou vẫn là Ojou thôi.}''

\chrint

Onidoro Ayame (\ruby{鬼}{おに}\ruby{灯}{どう}\ruby{彩}{あや}\ruby{芽}{め}, \ruby{Ô-ni}{Quỷ}-\ruby{đô-rô}{Đinh}\ \ruby{A-gia}{Thái}-\ruby{Mê}{nha}) là chủ tịch hội học sinh kiêm thánh nữ nghịch ngợm của trường Aogami. Dưới mái tóc trắng với vài viền đỏ lúc nào cũng tết bím đung đưa, cô cai quản hội đồng bằng ``bàn tay sắt bọc trong trò đùa'': buổi sáng có thể bất ngờ đổi nhãn tủ đồ, buổi trưa phát thanh giả mạo thông báo ``hôm nay học sinh được nghỉ tiếp'', chiều thì treo poster sự kiện ``thi ăn bánh crepe tốc độ'' chẳng ai hiểu nguồn gốc. Những màn này khiến học sinh vừa cười vừa rên, nhưng không ai dám phản kháng, bởi bất cứ ai vi phạm nội quy đều sẽ ``tình cờ'' bị gọi lên văn phòng hội học sinh, nơi Ayame chờ sẵn với bảng điểm, ảnh chứng cứ và nụ cười thánh thiện đầy ám khí.

Đằng sau vẻ nhây nhây ấy lại là bộ óc siêu tổ chức: lịch học, lịch thi, lễ hội văn nghệ, hội thao\dots\ tất cả được cô lập trình chặt chẽ đến từng phút, khiến giáo viên thở phào và ban giám hiệu không tiếc lời khen. Học sinh kính nể, kẻ phá phách thì khiếp vía; tiếng ``Hội trưởng hội học sinh Ayame có việc muốn nói'' đủ để bất kỳ đứa cá biệt nào chỉnh tề cà vạt trong vòng ba giây.

Thỉnh thoảng, những trò đùa của cô lại mang mùi\dots\ khảo nghiệm. Có lần cô bày trò ``đêm săn đèn lồng'' trong khuôn viên tối om, dẫn dắt học sinh đi qua các hành lang vắng; lần khác thì phát tờ rơi mời ``tham quan tầng hầm cũ chưa ai mở cửa''. Người ta cười đùa tham gia, nhưng ánh mắt Ayame lúc ấy lấp lánh không chỉ tinh quái, mà còn như đang thăm dò xem liệu có ai cảm nhận được ``điều gì đó'' đang lẩn khuất dưới lớp sơn mới của ngôi trường bình yên. Có đứa bảo cô chỉ trêu ghẹo; có đứa thì thầm rằng, đằng sau tiếng cười, Ayame đang dò đường cho cả hai phe: kẻ bình thường và thứ không bình thường đang ngủ quên trong bóng tối.

\chrint

``Được rồi Onidoro-san, và\dots\ người cuối cùng. Em lên đi.''

Cô gái tóc vàng dịu dàng cùng mái tóc bồng bềnh dài tới sống lưng tiến lên phía trước, bế theo một hộp bánh quy.

``Chào mọi người\dots\ mình là Hitsujikado Wata, phó hội trưởng hội học sinh. Mình thích nướng bánh, và\dots\ nếu mọi người muốn thử thì\dots\ ừm\dots\ mình nghĩ cũng được.''

Một vài tiếng ``dễ thương quá!'' bật lên từ lớp.

Kuon nhìn cô mà thấy lòng ấm:
``\textit{Làm mình liên tưởng tới Wata-chan. Hiền, mềm, êm như một bài hát ru.}''

\chrint

Hitsujikado Wata (\ruby{羊}{ひつじ}\ruby{角}{かど}\ruby{綿}{わた}, \ruby{Hi-xư-gi}{Dương}-\ruby{ca-đô}{Giác}\ \ruby{Oa-ta}{Miên}), một thiếu nữ hiền hòa với mái tóc bồng bềnh mềm mại khiến người đối diện chỉ muốn gục vào mà nghỉ ngơi.

Từ thuở ấu thơ, cô đã là bạn thân thiết của Ayame; khi người bạn kia sắp vượt khỏi giới hạn trò đùa, Wata lập tức trổ tài ``phó tướng'' đầy lý trí, kéo cả hai trở lại mặt đất. Với vai trò phó chủ tịch hội học sinh, cô không chỉ giữ trật tự mà còn giữ lấy niềm tin của thầy cô và bạn bè.

Tính cách bình thản, lời nói chậm rãi nhưng đầy thấu hiểu biến cô thành ``chiếc hộp thư sống'' của lớp, nơi mọi người có thể gõ cửa giữa những buổi chiều căng thẳng. Sau mỗi cuộc họp câu lạc bộ, căn phòng nhỏ lại ngập hương bơ sữa ấm; Wata mang đến những chiếc bánh quy vừa nướng, chia đều cho từng người như cách cô chia nhỏ nỗi lo trong lòng họ.

Dưới vẻ ngoài ấy là một sức mạnh lặng lẽ được rèn từ những lần Ayame gục ngã, cô đã học cách đứng vững để làm chỗ dựa, để khi thế giới rung chuyển, cô vẫn giữ được nhịp thở đều đặn cho cả hai.

\chrint

Sau màn tự giới thiệu của cả sáu, Yuko khẽ nâng khóe môi, giọng ấm áp phủ xuống lớp học.

``Cảm ơn các em đã chia sẻ. Lớp 1-C chúng ta có thêm những thành viên mới đầy sức sống rồi đây. Mong rằng tất cả sẽ cùng nhau học tập, trưởng thành và tạo nên những kỷ niệm đẹp trong năm học này.''

Koishin nhíu mày một chút, còn lại năm người đồng thanh cúi chào; tiếng họ nhẹ nhàng nhưng đủ vang khắp bốn bức tường: ``Xin cả lớp hãy giúp đỡ bọn mình trong thời gian sắp tới!''

Phòng học vỡ òa tiếng vỗ tay, rộn rã như hội chợ đầu năm.

Honoka hai mắt sáng lấp lánh, hơi nhún gót: ``Waa\~{} cuối cùng cũng có thêm bạn mới! Cậu nào muốn ngồi cạnh mình không? Mình có kẹo dẻo nè\~{}!''

Mitsuki đặt nhẹ hai ngón tay dưới cằm, khẽ gật: ``Chào mọi người. Nếu cần tài liệu ôn thi hay chỉ muốn tìm nơi yên tĩnh đọc sách, cứ ghé bàn mình nhé.''

Saya chậm rãi vén tóc sau tai, giọng ngân như ru: ``Hể\~{} mình thích ngắm mây trên sân thượng\dots\ Nếu ai đó cũng thích, có thể ngồi cùng mình nha\~{}''

Mari xoay nhẹ bím tóc, cười tươi như bông hoa đầu hạ: ``Mình mang theo bút màu nước mới mua đây! Buổi ra chơi nào rảnh, mình sẽ vẽ tặng mọi người một bức nhé!''

Saki khẽ hít hà, ôm cánh tay mình: ``Hương hoa hồng thật dễ chịu\dots\ À, mình có thể giúp cắm bình hoa nhỏ cho lớp nếu ai muốn\~{}''

Suzuki cúi đầu thật thấy, giọng mềm như lụa: ``Bọn mình sẽ giúp đỡ các cậu thật nhiệt tình! Onii-sama và Onee-sama thì sao ạ?''

Kaigou khúc khích, khẽ huých khuỷu tay vào sườn Kuon, thì thầm đủ cho cậu nghe: ``Kuon-kun, cậu thấy không? Lớp mình như bể cá mới, đầy màu sắc quá đi!''

Kuon nhìn lên trần, môi mấp máy thật khẽ: ``\hl{\dots Đúng vậy. Dù đổi sang thế giới này, mình vẫn ngửi được mùi\dots{} quen thuộc ấy.}''

\ct

Ngay lập tức, không khí lớp học náo nhiệt hẳn lên. Tiếng xì xào lan nhanh từng dãy bàn, ánh mắt tò mò bám theo bước các cô gái mới.

Ánh nắng sớng hắt vàng lên tóc Ayame đang cười toe toét, còn Wata khẽ cúi, tay nắm chặt sợi dây cặp như muốn trốn vào bóng của chính mình.

Hotaru khoanh tay hỏi thẳng: ``Ủa, hai cậu là hội trưởng với phó hội trưởng học sinh từ trước luôn á? Nhưng hôm nay mới nhập lớp mới mà?''

Ayame chống nạnh: ``Thì bọn tôi được bầu từ đợt khai giảng rồi. Lớp 1-C chỉ cập nhật chậm thôi. Hồi cấp hai tôi cũng là hội trưởng hội học sinh á\~{}'' cô nháy mắt như đang châm chọc cô nàng da ngăm.

Wata thì nhỏ nhẹ: ``Mình cũng thế. Chính vì có kinh nghiệm như vậy nên bọn mình giữ chức đó từ khi được xếp lớp luôn.''

Nanase bước đến, ánh mắt sắc bén kiểu ban kỷ luật: ``Tức là hai cậu sẽ cùng tôi xử lý vi phạm trong trường sao?''

Ayame cười lém lỉnh: ``Xử lý? Không, không! Tôi thì lo quan hệ đối ngoại là chính\~{} Còn Wata-chan mới là người ghi biên bản.''

Wata gật đầu: ``Mình ghi đẹp nắn nót lắm đấy.''

Nàng tóc xanh thở dài: ``Thôi\dots\ miễn sao làm đúng quy định. Nhưng mà đừng có giở trò gì đấy. Tôi sẽ để mắt tới hai cô.''

Ayame vỗ vai cô: ``Sát cánh nha, đồng nghiệp\~{}''

Ở một góc khác của lớp học, bộ ba Azuki-Hanabi-Inari đang tụm năm tụm ba bên chiếc ghế cuối cùng gần tủ đồ. Không khí có phần sôi động hơn hẳn, bởi Hanabi, đang tíu tít như con sóc quay tít.

Cô bất ngờ chạy vòng qua hai dãy bàn, mái tóc búng ra sau lưng như ngọn lửa nhỏ, rồi dừng lại trước mặt Azuki, người vừa mới cất chiếc hộp đựng bento vào ngăn bàn.

Cô bạn tóc nâu hơi khom người, hai tay chống nạnh, mắt sáng như đèn pha ô-tô: ``Này này, Azuki-chan! Cậu hát thử một đoạn được không!? Một đoạn thôi cũng được!''

Cô nàng tóc bob nghe vậy, gãi gãi đôi má đỏ ửng của mình: ``M-Mình cũng không chắc nữa\dots''

Inari vội chạy theo kịp Hanabi, khẽ kéo vạt áo blazer của cô nàng lại, nhỏ giọng nhắc nhở: ``Hanabi, đừng có làm người ta sợ\dots\ Azuki-chan mới chuyển trường mà.''

Cô nàng nghe vậy mới hạ nửa cái cằm xuống, nhưng đôi mắt vẫn long lanh không kém. Cô quay sang Azuki, hạ thấp giọng xuống một nấc, gần như cầu khẩn: ``Chỉ cần một đoạn nhỏ, mình hứa sẽ không ghi âm lén, cũng không réo lên như vừa nãy nữa!''

Azuki cúi xuống, hai ngón tay đan vào nhau trước bụng. Mái tóc nâu đỏ của cô khẽ rung khi cô thở ra một hơi dài, rồi bất ngờ nở nụ cười nhẹ, nụ cười khiến bầu không khí xung quanh như có thêm mùi hoa anh đào sớm.

``N-Nếu cậu thực sự muốn nghe thì sau giờ học, ở phòng nhạc trên tầng bốn, mình có thể hát thử một bài. Nhưng cậu phải hứa không được cười nha.''

``Yatta—!!''

Hanabi suýt bật nhảy lên, hai tay giơ cao như vừa ghi bàn ở phút 90. Inari thở phào, buông tay khỏi vạt áo Hanabi, mỉm cười an ủi:

``Cảm ơn cậu, Azuki-chan. Mình cũng muốn được nghe\dots\ nhưng mình sẽ ngồi xa một chút, đề phòng Hanabi quá khích.''

``Này! Mình đâu có đến mức đó!'' Hanabi phồng má. Và rồi, ba cô gái cùng cười, tiếng cười nhỏ nhẹ nhưng đủ khiến góc lớp cuối kia trở nên ấm áp hẳn.

Ở một diễn biến khác, Shion bất ngờ chạy tới, nắm cổ tay Shino kéo cô lại gần đến mức hai trán suýt chạm nhau. Đôi mắt tím biếc của nàng ``ảo thuật gia tự xưng'' sáng lên như phát hiện kho báu: ``Ngươi! Ta cảm nhận được khí chất huyền bí quanh ngươi! Hơi thở của bóng tối, tiếng thì thầm của sao chổi, tất cả đang hòa quyện trong từng sợi tóc ngươi! Ngươi có muốn gia nhập CLB Huyền Bí không!?''

Shino đỏ bừng mặt, hai tay vung vẩy như người chới với giữa dòng nước. Cô cười gượng, giọng lí nhí: ``À\dots\ ừm\dots\ cảm ơn lời mời của cậu, nhưng\dots\ tớ chỉ giỏi ở nhà bếp thôi. Cái khí chất huyền bí đó chắc là tớ không rành đâu.''

Touka đứng sau, xoa trán, thở dài như người đã quá quen với những màn ``tuyển quân'' đột xuất kia: ``Shion-san, làm ơn tha cho cậu ấy đi. Không phải ai cũng hiểu được mấy cái thần linh ma quỷ mà cậu nói đâu. Cô ấy đã run như cầy sấy rồi kìa.''

Đối diện họ, Saya đang bước đi về bàn như người mộng du. Bất chợt, cây bút máy trượt khỏi tay, rơi xuống nền gạch.

*CẠCH*

Koishin nghe vậy, lập tức cúi xuống nhặt trong nháy mắt, nhanh đến mức tóc khẽ bay. Cô đứng dậy, mặt nghiêng nghiêng, đưa bút về phía Saya, giọng nhỏ như tiếng lá rơi: ``Của cậu\dots\ này.''

Mặt cô đỏ bừng, lan xuống tận cổ áo sơ mi.

Nàng tóc vàng ngắn chớp mắt, ngạc nhiên xen lẫn cảm kích:
``Cảm ơn cậu nhiều\~{}''

``Hmph!'' Koishin quay phắt đi, tai đỏ như cà chua vừa nhúng nước sôi. ``Tại cậu vụng quá thôi, không phải tôi giúp đâu! Đừng\dots\ đừng hiểu lầm!'' Cô bước nhanh về phía trước, nhưng bước chân loạng choạng suýt đâm vào tường, để lại Saya đứng chôn chân, nở một nụ cười ẩn chứa muôn vàn tinh tú.

Ở phía bàn trên cùng, Aku và Kuon nhìn nhau lần nữa, ánh mắt như muốn nói ``chúng ta lại gặp nhau rồi''.

Cô nàng tóc hồng-xanh lí nhí, hai tay véo nhau trước bụng, mặt đỏ ửng: ``Cảm ơn vì lúc sáng đã đỡ tớ\dots\ tớ xin lỗi\dots\ tớ không nhìn đường\dots\ Cậu có bị thương không?''

Kaigou khoanh tay, mắt sáng quắc, giọng đầy nghi ngờ: ``Khoan, lúc sáng là sao? Kuon-kun lại gây chuyện gì đấy?''

Suzuki nghiêng đầu, mái tóc ngắn bay nhẹ, nhìn Kuon như thể vừa phát hiện ra một bí mật:
``Onii-sama, anh quen bạn ấy à? Sao em không nghe anh kể bao giờ?''

Kuon thở dài, xoa trán, rồi kể lại chuyện Aku suýt ngã lúc trên đường đến trường, cách cô lao như tên bắn, cách cậu đỡ được vai cô trước khi cô ``ăn mặt bàn'', và cách cô biến mất chỉ sau một câu cảm ơn lắp bắp. ``Chuyện chỉ có như vậy thôi.''

Kaigou và Suzuki đồng loạt kéo dài ``Àaaaa'', rồi liếc Kuon đầy\dots\ nghi vấn khó hiểu, kiểu ``Onii-sama/anh đang giấu chuyện gì đúng không?''.

Cậu Tổng thở dài một hơi nữa. ``\textit{Cơ mà nếu mà nói mình quen cô ấy thì\dots\ cũng đúng một phần.}''

Cả phòng học trở nên sống động hơn bao giờ hết: nhóm bên trái thử bánh quy Wata, nhóm bên phải xin Shion đọc thần chú ``làm bánh quy không béo'', nhóm giữa thì bàn chuyện câu lạc bộ chiều nay. Tiếng cười, tiếng ``wow'', tiếng ``có thật không?'' hòa vào nhau như bản giao hưởng học đường.

Trong khi đó, Kuon nhìn quanh, mồ hôi lạnh chảy dọc sống lưng, và nhận ra một sự thật tréo ngoe: \emph{cậu chính thức là con trai duy nhất của cả lớp 1-C—một ốc đảo nam tính giữa biển nữ sinh.}

``\textit{Giống hệt một bộ rom-com học đường harem đời thực\dots\ Ai sắp xếp cái số phận này vậy trời\dots}''

Cậu chỉ biết ngồi xuống, thở nhẹ như người vừa nhận án tù chung thân mang tên ``Harem Route Unlocked''.

\ct

Tiết chuẩn bị bài học của lớp tiếp tục với nhịp độ thoải mái mà Tsukishita Yuko đã duy trì từ đầu. Cô bước xuống bục giảng, trên tay là một xấp giấy dày, rồi nhẹ nhàng phát từng tờ cho từng bàn.

``Đây là phiếu đăng ký câu lạc bộ. Hạn nộp là ba ngày nữa. Các em nên suy nghĩ thật kỹ, vì hoạt động câu lạc bộ sẽ chiếm khá nhiều thời gian sau giờ học.'' Giọng của cô mềm mỏng, nhưng vẫn đủ cứng rắn để khiến cả lớp cảm nhận được tính nghiêm túc của vấn đề.

Ngay khi Yuko vừa quay trở lại bàn giáo viên, cả lớp 1-C lập tức nhao nhao như một phiên chợ nhỏ, bàn bạc xem câu lạc bộ nào sẽ phù hợp với sở thích của họ.

Miku đặt tờ đơn xuống bàn, gõ nhẹ cây bút lên mép giấy. ``Cậu định vào CLB nào, Mitsuki-san?''

Mitsuki chống cằm, ánh mắt mơ màng như đang nhìn một điểm xa xăm mà chẳng ai thấy.
``Mình cũng hơm biết nữa\~{} Cậu thì sao, Miku-chan?''

Cô nàng bờm Sói đen mỉm cười, hai má hây hây đỏ: ``Tôi ngưỡng mộ những người phụ nữ trưởng thành lắm, nên tôi đã đọc rất nhiều tạp chí thời trang. Cậu biết không, tôi cảm thấy thời trang như một thứ ngôn ngữ kiêu kỳ, chỉ cần nhìn là hiểu được tâm hồn người mặc.''

Mitsuki mở to mắt, khe khẽ: ``Ồ\dots\ tinh thế ghê ha\~{}''

Miku vội cúi xuống, day day mép tờ đơn: ``À không\dots\ chỉ\dots\ thích thôi mà.''

Mitsuki nghiêng đầu, nhìn quanh dãy hành lang rồi mỉm cười: ``Mình thường lang thang ở khu mua sắm phía đông thị trấn Aogami lúc tan học. Có một tiệm nhỏ treo mấy bộ váy thiết kế lạ mắt lắm, toàn là họa tiết đan xen pha nét cổ điển và hiện đại. Mình đứng ngắm mãi, thậm chí còn về nhà phác thảo lại.''

Cô khẽ vỗ vai Miku, giọng khích lệ: ``Nếu chỉ 'thích thôi' thì đã đủ để cậu bắt đầu rồi. Cảm hứng ban đầu luôn từ một chút yêu thích như vậy đấy.''

Miku cười, gật đầu: ``Vậy\dots\ chúng ta vào CLB Thiết kế thời trang nhé?''

``Ừ!'' Miku ngẩng lên, mắt sáng rực. ``Cùng nhau thôi!''

Trong khi đó, Kana đang hí hoáy viết, Kanade ngó sang. ``Bà chọn CLB gì vậy?''

``CLB Phát thanh!'' Kana trả lời nhanh như bắn súng, mắt sáng lên. ``Tôi muốn được nói thật nhiều, được người khác nghe thấy giọng mình. Biết đâu sau này thành DJ nổi tiếng thì sao?''

Kanade khẽ cười, véo nhẹ má bạn: ``Nghe đã nhỉ. Vậy tôi cũng vào luôn, làm kỹ thuật viên âm thanh. Đằng nào tôi cũng thích ngồi một mình trong phòng thu, lắng nghe nhịp sóng.''

Nàng tóc tím trợn mắt, véo lại: ``Ha?! Lý do đó cũng được tính à?''

``Hợp lý mà.''

Đối diện với họ, Yae nâng tờ đơn lên ánh sáng, rồi chép miệng, nhíu mày.

``Tớ đang phân vân quá\dots\ Mỹ thuật thì vẽ được, Thể thao thì chạy khỏe, nhưng chỉ được chọn một thôi,'' cô thở dài, lắc đầu. ``Thiếu cái gì đó\dots\ như thể cả hai đều chỉ là nửa vậy.''

Sakura liếc sang, giọng đều đều nhưng ẩn chút háo hức: ``Tôi thì chọn CLB Huyền bí.''

Yae suýt đánh rơi tờ đơn, mắt tròn xoe:``H-Hả? Tại sao vậy?''

``Thích thôi.'' Sakura nói tỉnh bơ, khẽ đẩy gọng kính. ``Những điều kỳ quái\dots\ đôi khi rất thú vị. Và có thể, sẽ giúp tôi hiểu thêm chút về bản thân thì sao?''

Yae thở dài, gật gù: ``Ừ thì\dots\ cậu hợp với nơi đó thật.'' cô quay lại nhìn tờ đơn, môi mím chặt. ``Còn tớ\dots\ chắc sẽ chọn Thể thao. Mỹ thuật thì thích thật đấy, nhưng mình lại có duyên với những hoạt động đồng đội hơn. Không muốn sau này hối tiếc vì mình chưa thử hết sức.''

Bốn tờ giấy được trượt nhẹ lên bàn giáo viên. Giáo viên chủ nhiệm Yuko gật đầu, ghi tên, rồi mỉm cười bí hiểm: ``Hy vọng các em sẽ tìm thấy điều mình cần ở nơi mình chọn. Và nhớ\dots\ đôi khi, CLB không chỉ là CLB, mà còn là cánh cửa.''

Cả bốn đứa nhìn nhau, không ai hiểu hết ý thầy, nhưng trong lòng đều thấy rộn ràng một thứ gì đó\dots\ như tiếng sóng vỗ nhẹ, báo hiệu mùa mới đã bắt đầu.

Ở quanh bàn Hikawa, Shino nhíu mày nhìn chiếc phiếu, đầu ngón tay gõ gõ nhịp nhàng lên mặt bàn. ``Hội học sinh\dots\ nghe có vẻ đáng sợ thật\dots\ nhưng chắc là mình sẽ thử.''

Cô thở dài một hơi nhỏ, ánh mắt thoáng xa xăm. Sau mớ vòng lặp hỗn loạn kia, cô hiểu rõ mình cần gì nhất: không phải là thêm thuật toán, mà là thêm can đảm. ``Mình phải tập đứng trước đám đông, tập nói mà không run\dots\ nếu không, mình mãi chỉ là con nhái trong lọ.''

Suzuki ngồi đối diện, hai tay ôm bình nước, mắt lim dim như mèo. ``Mình thì\dots\ chắc sẽ thử CLB Văn học.''

Shino ngẩng phắt lên: ``Hả? Cậu không vào Khoa học nữa à?''

``Ừ.'' Suzuki gãi má, ngại ngùng. ``Hồi trước, mình từng đọc trộm thư viện của bố, cả tủ sách cổ điển lẫn light-novel. Đôi lúc mình cũng chìm vào thế giới của những cuốn sách đấy\~{}'' cô cười nhẹ, má lúm đồng tiền ẩn hiện. ``Không bù cho cậu, dũng cảm như chuẩn bị leo núi.''

Shino bật cười, cái vỗ bàn trở thành vỗ vai. ``Được rồi, mỗi đứa một phía. Nhưng mà\dots\ nếu tớ có run quá, liệu cậu vẫn sẽ giúp tớ chứ?''

``Dĩ nhiên rồi!'' Suzuki đáp, đôi mắt cong lên đầy vui vẻ. ``Nhưng mình tin là cậu có thể làm được mà không cần mình đâu. Cứ tự tin vào bản thân là được rồi.''

Khắp phòng học, từng nhóm bạn ép nhau điền vào phiếu. Sau một hồi bàn bạc, thuyết phục, thậm chí là chèo kéo, cuối cùng tất cả các cô gái của lớp 1-C đã chọn được câu lạc bộ.

Không khí trong lớp sôi nổi đến mức Yuko nhìn từ xa mà bật cười nhỏ. ``\textit{Cứ như chợ phiên vậy\dots\ nhưng mà, kể ra cũng hay.}''

Chỉ còn lại Kuon, cậu con trai duy nhất của lớp. ``Kuon-kun, cậu định gia nhập vào câu lạc bộ nào vậy?''

Kaigou liếc sang tờ đơn của cậu Tổng, rồi lập tức\dots\ bật lùi khỏi ghế. ``Ể!??''

Suzuki ghé sang xem từ tiếng thét của cô chị, và rồi đơ người toàn tập. ``O-Onii-sama\dots''

Trên tờ đơn của cậu Tổng, chỉ có đúng ba chữ: 「\ruby{帰}{き}\ruby{宅}{たく}\ruby{部}{ぶ}」 (Câu lạc bộ Về nhà).

Cả lớp nghe Kaigou la lên thì đồng loạt quay sang.

Một giây.

Hai giây.

Ba giây.

``\textbf{ỂỂỂỂỂ!?}''

Tiếng bàn tán bùng nổ như ong vỡ tổ. Có người trố mắt, có người che miệng, có người đập bàn đứng phắt dậy.

Uzuki trố mắt: ``Hả!? Kuon-san không vào câu lạc bộ nào á? Tôi tưởng cậu có khả năng chụp ảnh lắm chứ?''

Shiina gõ gõ trán, cười như thể vừa thấy quái vật lạ: ``Đỉnh thật. Lười mà còn hợp pháp nữa. Tôi đang tính rủ cậu vào bộ môn Karate, chứ còn gì!''

Sakura thì chỉ ``Ồ\dots'' đúng một chữ, nhưng đôi mắt khẽ liếc sang phía cậu Tổng, ánh lên vẻ ``tôi đã biết từ đầu rồi''.

Một nhóm phía sau còn xì xầm to hơn:

``Cậu ta không sợ bị điểm trừ sao?''

``Nghe đồn hiệu trưởng rất khắt khe về chuyện hoạt động ngoại khóa mà\dots''

``Hay là cậu ấy có bí kíp gì để thoát được?''

Ở phía bục giảng, Yuko liếc lên, mỉm cười nửa như hiểu chuyện, nửa như đang ghi chú vào trí nhớ: ``\textit{Loại học sinh chọn `về nhà' cũng không phải hiếm. Nhưng mà tại sao em ấy lại quyết định như vậy nhỉ?}''

Dưới lớp, những lời bàn tán xôn xao vẫn không ngừng.

``Kuon-san!? Sao cậu lại chọn cái này?'' Shino hỏi thẳng.

Kuon thản nhiên đáp lại: ``Tôi không hợp thể thao, không thích hoạt động tập thể, không có sở thích rõ ràng nào để gắn bó. Chưa kể còn phí thời gian. Về nhà là lựa chọn hợp lý nhất.''

Mitsuki nghiêng đầu: ``Ngắn gọn thật đó. Lại còn thẳng đến phũ phàng nữa chứ.''

``Nhưng ít nhất thì cậu cũng nên--'' Uzuki định chen vào thuyết phục, nhưng cậu Tổng chỉ giơ tay lên như ra dấu dừng lại.

``Không cần phải lo cho tôi,'' cậu đáp tỉnh queo, không một chút dao động.

Đúng lúc đó, Koishin khẽ hừ một tiếng, đôi mắt cô liếc xéo, giọng đầy khinh khỉnh: ``Hmph. Đúng kiểu lười. Con trai gì mà yếu đuối thế.''

Không khí trong lớp như chùng xuống một nhịp, mọi người đồng loạt khựng lại, mắt liếc nhau.

Kuon không thèm ngoảnh đầu, chỉ nhàn nhạt đáp, giọng lạnh như dao cắt giữa trưa hè: ``Không cần cậu phải xía vào.''

Một câu ngắn gọn, không chút cảm xúc, nhưng đủ khiến không gian thêm cứng lại.

Shiina lập tức bước chen vào, tay chống nạnh, mắt trợn trừng: ``Ê! Đừng xúc phạm đàn ông duy nhất của lớp chứ!''

Nhưng cô nàng tóc vàng ruy băng đỏ chỉ bĩu môi, quay mặt đi, không thèm đáp lại.

Kaigou nhỏ giọng lầm bầm: ``\hl{Kuon-kun thản nhiên đến đáng sợ\dots{} y hệt như một tổng quản\dots}''

Yuko vỗ tay ba cái, nhẹ mà vang. ``Rồi, thảo luận nhiêu đó là đủ rồi. Cất phiếu vào cặp, lát nữa tan học nhớ nộp nhé. Giờ vào bài mới.''

Cả lớp nhanh chóng ổn định lại, tiếng ghế kéo nhẹ hòa cùng tiếng lật sách.

Lớp 1-C chính thức bước vào tiết học đầu tiên sau tiết truy bài đầu giờ.

\ct

\pov{Hikawa Kuon}

Tiết Sinh học đầu tiên của năm học. Tôi vốn không kỳ vọng gì nhiều, vì vốn dĩ môn Sinh học không phải là sở trường của tôi (ý tôi là, tại sao tôi phải để ý một môn mà chỉ việc học thuộc lòng thôi chứ?)

Nhưng rồi, cách mà Yuko-sensei dạy học khác với tưởng tượng của tôi. Cô đứng trước bảng, cầm phấn bằng tay phải, tay trái đặt hờ lên mép bàn giáo viên. Giọng cô trầm, mềm nhưng vẫn đủ để khiến cả lớp im lặng khi bắt đầu bài giảng.

``--Và đây là quy trình quang hợp. Các em lưu ý rằng các sắc tố quang hợp không chỉ có mỗi chlorophyll\footnote{Sắc tố màu xanh lá cây có trong thực vật, chứa ít nhất một nguyên tử Magnesi (magiê). Còn được biết với cái tên là `diệp lục'.}.''

Sakura-san ở bàn ngẩng lên ngay lập tức, ánh mắt sáng như vừa nghe thứ gì mình yêu thích nhất trên đời. Thứ này đúng là sở trường của cô ấy thật, quả không hổ danh nhà môi trường học có khác.

Tôi chống cằm, nhìn cảnh cả lớp ghi chép lạch cạch. Lớp 1-C im lặng hơn tôi tưởng, ít nhất là cho đến khi\dots

*BỘP*

Một tiếng động sắc gọn vang lên như một phát đạn nhỏ. Tôi liếc sang. Bên cạnh cửa sổ, Yuka-san đang gục trên bàn ngủ ngon lành, còn Akane-san ngồi bên cạnh thì cố lay vai cô ấy.

``\hl{Yuka! Dậy, dậy mau!}''

Nhưng Yuka hoàn toàn không phản ứng. Tai thì cụp xuống, thỉnh thoảng giật giật.

Rồi tôi nghe tiếng gió xé nhẹ, hay đúng hơn là tiếng đồ vật xé gió. Một vật trắng bay thẳng theo đường cong hoàn hảo.

*CHÁT*

Viên phấn đập thẳng vào trán cô bạn bờm mèo đó.

``Ư—!'' Yuka-san bật dậy như lò xo, hai tay ôm đầu. ``Đau đau đau\dots! Ai đánh tôi vậy?!''

Yuko-sensei vẫn quay mặt lại bảng, bình thản đến mức tôi nghi ngờ cô có cả kỹ năng ngắm bắn không cần nhìn.

``Không ngủ trong tiết học của tôi,'' cô nói nhẹ như gió thoảng.

Cả lớp rủ rỉ cười. Tôi thầm nghĩ: ``\textit{cô ấy mà bị xử thế này, mấy đứa còn lại coi như xong.}''

\ct

Tiết học môn sinh cứ diễn ra bình thường như vậy, không có chuyện gì xảy ra.

Đó là cho tới khi gần cuối tiết. Khi Sakura-san hăng hái ghi chép về cấu trúc tế bào thực vật, tôi bất giác liếc sang cô ấy. Bút trên tay Sakura khựng lại. Ánh mắt cô hướng ra cửa sổ, hơi nheo lại, sắc dần, lạnh dần.

Tôi khẽ cau mày. ``\textit{Sakura-san\dots?}''

Cô không trả lời.

Ở hàng ghế kế bên, Yuka-san dù vừa bị đánh thức lại cau mày trong cơn ngái ngủ. Tai cô giật nhẹ, như đang phản ứng với thứ gì đó vô hình.

Tôi thoáng chột dạ. Năng lực ngoại cảm của cô ấy.

Không khí trong lớp đột nhiên khác hẳn.

Kaigou-chan đang vẽ đường thẳng trong vở thì dừng lại giữa chừng. Suzuki ngẩng đầu khỏi tập ghi chép, ánh mắt hơi co giật như bắt được tín hiệu lạ.

Shino thì nhìn chăm chăm ra cửa sổ từ bao giờ. Tôi chưa kịp hỏi thì cô ấy lẩm bẩm: ``\dots Có thứ gì đó đang đứng ngoài sân kìa.''

Tim tôi giật mạnh một nhịp. Tôi nhìn theo hướng của cô ấy\dots\ và thấy \emph{nó}.

Giữa sân trường, ngay chính giữa khoảng đất nơi học sinh thường xếp hàng buổi sáng, có một hình dáng nhỏ bé đứng lẻ loi. Một đứa trẻ mặc kimono cũ kỹ, tà áo rách nhẹ ở mép.  Nhưng điều đáng sợ hơn: tà áo khẽ đung đưa. Dù trời hoàn toàn không có gió. Khuôn mặt nó mờ như bị xoá nhoè. Không mắt. Không mũi. Không miệng. Chỉ là một khoảng trống lờ mờ hình người. Hoặc có lẽ, góc nhìn ở trên cao khiến tôi không thể nhìn thấy rõ. Thứ mà chúng tôi đã nhìn thấy vào chiều tối hôm qua.

Sân trường vắng tanh, ngoài mấy con chim sẻ đang nhặt hạt, nhưng đúng lúc tôi nhìn, đàn chim xoạt một tiếng, bay vọt lên tán loạn như vừa bị quấy động bởi thứ gì đó kinh khủng.

Tôi siết chặt bút. Cơ thể lạnh dần. Cảm giác như có luồng khí lạnh chạy dọc xương sống, bủa quanh ngực.

Từ cổ tay, giọng Game vang lên rất nhỏ, như thì thầm xuyên qua làn da: ``\eng{\hl{Là một trong số chúng.}}''

Tôi hít vào thật sâu nhưng ngực vẫn nặng trĩu. Không khí trong phòng học như đặc quánh lại. Và đứa trẻ ấy\dots\ \emph{đang đứng bất động nhìn về phía lớp chúng tôi}.

Tôi nheo mắt, cố quan sát kỹ hơn cái bóng nhỏ đang đứng giữa sân thì\dots

``\textbf{\dots Bố ơi\dots\ Mẹ ơi\dots\ Cứu con với\dots}''

\dots Tôi nghe thấy một giọng nói cao vút. Rất khẽ. Nhưng đủ để khiến tôi nổi da gà ngay lập tức. Chính xác hơn, là tiếng trẻ con khóc. Không rõ là bé trai hay bé gái, chỉ là thứ âm thanh nhỏ, yếu, như phát ra từ một nơi nào đó sâu thẳm dưới mặt đất. Tiếng khóc tắc nghẹn, đứt đoạn, rồi biến mất.

Tôi giật mình, quay sang Sakura-san,  cô ấy cũng đang run nhẹ đầu bút trong tay. Shino-san nghiêng đầu, ánh mắt sắc như muốn xuyên qua không gian. Yuka-san ôm lấy tai, môi run bần bật. Suzuki thì nắm chặt mép bàn tới trắng cả bàn tay. Kaigou-chan khựng lại, mắt mở lớn hơn bình thường. Dường như không phải chỉ mình tôi nghe thấy.

Ngay khoảnh khắc tiếp theo, một giọng người lớn trầm khàn, tuyệt vọng, như từ đáy lòng đất vọng lên\dots

``\textbf{Xin tha cho chúng nó\dots\ xin đừng\dots!}''

Âm lượng không lớn, nhưng nặng như một tảng đá đè thẳng lên lồng ngực.

Giọng nói ấy không phải tiếng gió. Không phải trí tưởng tượng. Nó chân thực một cách đáng sợ.

Nhưng chưa dừng lại ở đó.

Tôi hình dung được việc đứa trẻ đó bị kéo đi, bị đẩy xuống hố tối, bị lấp đất lên người\dots\ nó nhói lên trong đầu tôi như ký ức của một ai đó khác đang trộn lẫn vào tâm trí mình.

Một hình ảnh thoáng hiện ra: một bàn tay bé nhỏ bị túm lấy. Tiếng đất lạo xạo. Tiếng gào thét bị nghẹn lại giữa cổ họng.

Không lẽ nào\dots\ đó là một nghi thức cổ xưa mà Sakura-san từng nói đến trong vòng lặp trước?

Tất cả đều lặp lại.

Tôi nắm chặt mép bàn đến mức khớp tay tê đi.

Tôi liếc lại về phía đứa trẻ giữa sân. Không hề di chuyển, không hề quay đầu.

Và nó bắt đầu tan ra. Từng mảnh. Từng hạt sáng nhỏ thoát khỏi cơ thể nó, lơ lửng như bụi vỡ dưới nắng.

Chúng bay lên cao, rồi biến mất hoàn toàn, như bị kéo trở về một nơi mà chúng thuộc về.

Tôi thở ra, nhưng ngực vẫn nặng.

``\textit{Tại sao nó lại xuất hiện ở đây?}'' tôi lẩm bẩm. ``\textit{Mục đích của nó là gì?}''

Yuka-san lên tiếng, tay bóp lấy đầu như đang cố giải mã thứ gì đó rất mơ hồ: ``Tôi cảm thấy\dots\ chúng không muốn làm hại ai. Những đứa trẻ đó giống như đang bị mắc kẹt giữa ba lớp ký ức khác nhau\dots\ như bị chia đôi vậy.''

``Tách đôi ký ức?'' tôi nhíu mày.

``Ừm\dots\ không hoàn toàn là của chúng nữa,'' nàng bờm mèo nuốt nước bọt. ``Có thứ gì đó kéo chúng sang hướng khác.''

Tôi định hỏi thêm thì bất chợt\dots

``Mấy em đang nhìn gì ngoài cửa sổ vậy?''

Giọng Yuko-sensei cất lên, bình thản nhưng sắc như lưỡi dao. Tôi giật mình quay lại. Cô vẫn đang giảng bài, nhưng ánh mắt thì rõ ràng hướng về phía sáu người chúng tôi. Ánh nhìn ấy trông không phải cái nhìn trách móc học sinh không tập trung. Đó là cái nhìn của một người lớn đang biết chính xác rằng có gì đó bất thường đang diễn ra.

``Không có gì đâu ạ, sensei,'' tôi đáp, rồi cố nặn ra một câu nói dối. ``Em chỉ nhìn mấy con chim bên ngoài thôi ạ.''

Sensei thở dài. ``Tập trung vào nghe giảng đi nhé. Cả các em nữa.'' Rồi, cô quay lại với bảng, tiếp tục giảng như chuyện vừa rồi chưa từng xảy ra.

``\hl{\dots Tôi nghĩ chúng ta nên tập trung vào bài giảng trước,}'' Game nói. ``\hl{Tạm gác việc đó ra một bên đi đã.}''

Tôi gật đầu, mắt liếc về bảng, tay tiếp tục viết những lời cô giảng. Tuy nhiên, dù đã cố bình tĩnh, tôi không thể quên được khoảnh khắc ấy.

``\textit{Chỉ hy vọng là Yuko-sensei sẽ không vướng vào chuyện này. Mình không muốn mọi thứ phức tạp thêm nữa.}''

Nhưng sâu trong lòng, tôi biết.

``\textit{Không sớm thì muộn, cảm giác này nói cho mình thấy rằng\dots\ Mình sẽ phải đối mặt với Yuko-sensei.}''

Bởi vì cô ấy đã nhìn thấy. Cô ấy biết. Và cô ấy không phải người bình thường.

Tiết học tiếp tục như không có gì xảy ra.

Nhưng với tôi, và những người đã nhìn thấy hồn ma kia\dots\ mọi thứ đã bắt đầu thay đổi.

\ct

Bốn giờ chiều, tiết cuối cùng trong ngày cũng vừa kết thúc. Sau giờ tan học, lớp 1-C dần vơi người. Tiếng ghế kéo, tiếng túi sách lách cách vang ra từng đợt nhẹ. Tôi lặng lẽ thu dọn sách vở, xếp từng quyển vở vào cặp.

Trong đầu vẫn thoáng hiện lên khoảnh khắc cuối tiết Sinh học buổi sáng, khi Yuko-sensei nhìn thẳng vào tôi và nói: ``Hikawa Kuon-san, em có thể ở lại vài phút được chứ?'' Tôi đã hơi khựng lại, nhưng cố giữ bình tĩnh, đáp: ``Dạ được ạ. Cô muốn hỏi gì em sao?''

Ánh mắt của cô khi ấy hơi cong lên, trông vừa hiền vừa sắc: ``Cũng không có gì nghiêm trọng đâu. Tôi chỉ muốn hỏi vài chuyện thôi.'' Dứt lời,cô rời lớp trước, để lại cả đám 1-C nhìn tôi như thể tôi sắp bị mang đi xét xử.

Ngay khi cửa lớp đóng lại, Kaigou-chan lập tức nhào đến bàn tôi: ``Cậu làm gì để bị gọi lại thế!? Ta vừa mới vào trường được có hai ngày đấy!''

Tôi thở dài: ``Tôi có làm gì đâu mà bị gọi.''

Suzuki khoanh tay, nghiêng đầu: ``Em cũng không biết nữa\dots\ ánh mắt Yuko-sensei khi gọi anh nó hơi giống kiểu `cô biết em đang giấu cái gì đó đấy'. Onii-sama cẩn thận thật.''

``Nghiêm trọng thế cơ à?'' tôi nói đùa, rồi lại thấy mình chẳng đùa được.

Cả lớp nhìn tôi với vẻ nửa lo lắng nửa tò mò, riêng Yae-san thì thầm với Uzuki-san kiểu ``Kuon-san bị cô gọi kìa\dots\ chết rồi\dots'', còn Koishin-san thì cười khẩy như thể tôi quả nhiên là ``học trò rắc rối''. Tôi chẳng buồn đáp lại, chỉ lặng lẽ rời khỏi lớp học trước khi có ai đó kịp nói điều gì.

\ct

Khi tôi tới phòng giáo viên, Yuko-sensei đang đứng cạnh cửa sổ, ánh sáng chiều vàng phủ lên mái tóc vàng đỏ của cô khiến nó trông như đang cháy nhẹ.

Cô ngoảnh lại, mỉm cười: ``Đến rồi à, Kuon-san. Vào đi.''

Tôi kéo cửa bước vào. Không khí có mùi phấn và giấy mới, dịu nhẹ nhưng có chút gì đó lạnh lẽo. Trước mặt tôi, cô đặt xuống bàn một tập hồ sơ. Sau đó, cô rút ra một tờ duy nhất, tờ đơn CLB của tôi. Chính cái tờ mà tôi chỉ viết đúng ba chữ: 帰宅部 - Câu lạc bộ Về nhà. Cô nâng tờ giấy lên, như đang soi từng nét chữ: ``Trước hết thì\dots\ Tôi có điều này thắc mắc.'' Cô cúi đầu hỏi, giọng mềm như nước nhưng vẫn có chút nghiêng về sự thẩm định: ``Vì sao em lại chọn `CLB về nhà' vậy, Kuon-san?''

Không trách mắng gì. Chỉ là thật sự tò mò. Tôi thở ra một hơi, nói thật: ``Thực ra, em không có sở thích gì nổi trội, và cũng không giỏi thể thao. Hồi cấp hai, em thậm chí còn suýt bị đánh trượt chỉ vì chạy 1.500m không nổi.''

Tất nhiên đó là chuyện xảy ra ở đại học của thế giới gốc, nhưng tôi phải sửa như vậy để lý do nghe thuyết phục hơn. ``Nhưng vì có ý thức học tập nên vẫn được cho qua\dots\ đại loại là vậy ạ.''

``Vậy à\dots'' Yuko-sensei gật, ngón tay khẽ gõ lên mép bàn. Cô hỏi tiếp, giọng chậm rãi, như muốn thử tôi: ``Thế còn Hội học sinh?''

Tôi cười nhẹ, lắc đầu: ``Hội học sinh đang có nhiều nhân lực lắm rồi ạ. Ayame-san làm hội trưởng, Wata-san làm phó, Nanase-san thì ở ban kỷ luật. Em cũng nghe Shino-san với Inari-san chuẩn bị vào nữa. Em mà vào thì cũng chẳng giúp được gì nhiều. Ngồi chơi không thì kỳ lắm\dots\ nên tốt nhất là đừng vào từ đầu.''

Ngụ ý của tôi - vào CLB chỉ để tốn thời gian - chắc là cô hiểu. Không cần phải nói thẳng ra, chỉ cần hiểu ý đằng sau là được.

Sau khi nghe những gì tôi trình bày, cô giáo nghiêng đầu, mắt hơi nheo lại như đang nghiên cứu biểu cảm của tôi.

Tôi bèn hỏi thẳng: ``Cô cho em hỏi\dots\ việc tham gia CLB có bắt buộc không ạ?''

Một cái ``Không,'' đầy dứt khoát, gọn lẹ.  ``Thế nên\dots'' tôi gật đầu, ``em mới chọn không vào CLB nào cả.''

Một khoảnh khắc im lặng. Yuko-sensei nhìn tôi, đôi mắt xanh nước biển trông mềm nhưng ẩn bên dưới là lớp phân tích sắc bén như dao. Cô bèn chuyển chủ đề: ``Còn một chuyện nữa. Em hỏi tôi trong giờ ăn trưa rằng học sinh có được đi làm thêm không. Vậy, em hỏi vậy để làm gì?''

Tôi đan tay lại, nói thật: ``Nhà em chỉ có ba anh em sống với nhau. Bố mẹ thì hầu như không có nhà, chỉ để lại tiền đủ ăn vài tuần.''

``Ừm\dots'' cô hơi cau mày, cảm xúc thoáng thay đổi.

``Thế nên là,'' tôi tiếp tục, ``em đang định tìm một công việc bán thời gian, để dành phòng khi cần thiết.'' Từ những kinh nghiệm tôi đã gặt hái được khi là một tổng quản, tôi biết rất rõ rằng của để dành là một thứ gì đó không bao giờ thừa, đặc biệt khi cuộc sống luôn ẩn chứa những biến cố bất ngờ.

Yuko-sensei đặt tờ đơn CLB xuống, đan tay trước ngực: ``\dots Và đó là lý do em không theo câu lạc bộ nào?''

Tôi gật: ``Đó cũng là một trong các lý do ạ. Mong cô thông cảm cho.''

Cô nhìn tôi thêm vài giây, đôi mắt xanh ấy lại có màu phân tích quen thuộc. Có chút tò mò. Có chút dò xét. Và có chút gì đó\dots\ như đang đánh giá xem tôi là kiểu học sinh nào. Kiểu dạng không phải là một học sinh nào mà cô biết. Không phải kiểu soi xét ác ý, mà là kiểu ``mình cần hiểu rõ đứa trẻ này để quản lý nó cho đúng''. Tôi cảm nhận rõ điều đó, và cảm giác nhẹ râm ran chạy dọc sống lưng.

Tôi khẽ nghĩ: ``Yuko-sensei\dots\ không đơn giản chỉ là giáo viên bình thường.''

Cô đang nhìn tôi như thể tôi có thứ gì đó khác biệt. Và tôi cũng không chắc cô sẽ để chuyện này yên đâu. Cô bất chợt đặt cây bút xuống, khoanh tay lại, nghiêng đầu như đang quan sát một mẫu vật sống trong phòng thí nghiệm. Rồi, bất ngờ, giọng cô hạ thấp hẳn xuống: ``Trong lúc dạy, tôi để ý một số em nhìn ra cửa sổ\dots\ hệt như nhìn thấy thứ gì đó. Nhất là em đấy, Kuon-san.''

Tôi khựng lại. Không trả lời. Trong đầu chỉ bật ra một ý nghĩ duy nhất: ``\textit{Mình không thể ngờ là Yuko-sensei nắm thóp được bọn mình\dots\ Quả đúng là cô giáo sắc sảo thật.}''

``Tôi là giáo viên Sinh học,'' Yuko-sensei tiếp tục, giọng bình thản đến khó đoán. ``Tôi biết rõ cơ chế phản xạ khi một người nhìn thấy. Một sinh vật, hoặc một mối đe dọa. Và phản ứng của sáu em - Kaigou-san, Suzuki-san, Misora-san, Nekomiya-san, Shinomiya-san, và em\dots\ không giống thứ gì tầm thường cả.''

Tôi nheo mắt nhẹ, quay mặt sang hướng cửa sổ, tránh ánh nhìn sâu như soi thấu nội tâm của cô. ``\dots Chuyện cũng rất phức tạp đấy ạ, thậm chí có gì đó hơi tâm linh. Nếu em nói, liệu cô có tin không?''

Yuko-sensei hơi mỉm, nụ cười không hẳn là vui, mà giống kiểu cô đã biết trước mọi câu hỏi tôi sắp nói.

``Tôi biết mọi thứ về thị trấn Aogami. Nơi này từng xảy ra sạt lở lớn. Tôi cũng biết về những báo cáo `tiếng khóc trẻ con' vào ban đêm.''

Cô đặt tay lên tập hồ sơ câu lạc bộ mà chúng tôi đã nộp ban nãy. Nhưng bên dưới tôi thấy thấp thoáng những tập báo cáo dày cộp với ký hiệu của đội khảo sát môi trường thị trấn.  ``Nơi này từng xảy ra sạt lở lớn. Đất bị xói mòn trắng, rồi\dots\ không hiểu vì sao lại tái sinh nhanh bất thường. Trong các báo cáo môi trường có ghi rõ. Và,'' cô lật nhẹ một trang tài liệu cũ, ``tôi cũng biết về những ghi nhận `tiếng khóc trẻ con' vào ban đêm. Không phải tin đồn. Đó là báo cáo dân sự mà đội sinh thái không giải thích được.''

Cô đặt ngón tay lên mép hồ sơ, suy tư giải thích: ``Khoa học không thể bác bỏ hoàn toàn. Dữ liệu đôi khi nói lên nhiều điều hơn những gì con người chịu tin. Và tôi tin vào đôi mắt học sinh mình hơn bất kỳ báo cáo nào.''

Tôi khựng lại. Không phải vì cô tin, mà vì cô biết quá nhiều. Hơn mức một giáo viên bình thường nên biết.

Cô tiếp lời, giọng thấp hơn: ``Phản ứng của mấy đứa lúc sáng\dots\ không phải phản xạ giật mình. Đó là kiểu phản ứng khi một người nhìn thấy\dots\ một sinh vật thật sự. Cơ mặt, đồng tử, nhịp thở\dots\ tất cả đều nói lên điều đó.''

Lồng ngực tôi siết lại. Tôi không nghĩ cô tin nhanh đến thế. Cũng không nghĩ cô quan sát kỹ đến mức nhận ra cả phản ứng sinh học của chúng tôi. Nhưng càng nhiều người biết, càng nguy hiểm. Tôi không thể kéo cô vào chuyện này, không một ai hết.

Tôi hít sâu một hơi, nói chậm rãi: ``Em xin lỗi, nhưng hiện tại\dots\ em không thể nói được. Hoặc ít nhất, không phải bây giờ.''

Đó là lời nói dối trắng. Nhưng tôi cần nó.

Yuko nhìn tôi một lúc lâu, rồi mới gật đầu. ``Được rồi. Tôi tôn trọng. Nhưng đừng để bản thân chịu đựng một mình.''

Tôi cúi đầu cảm ơn, đứng dậy.

``Vậy, nếu không còn việc gì khác\dots\ cho em xin phép ạ.''

Tôi vừa chạm tay vào tay nắm cửa thì Yuko gọi với ra: ``Mà này, Kuon-san\dots''

Tôi quay lại. ``Dạ?''

Cô nghiêng người ra phía ánh sáng từ cửa sổ, đôi mắt nghiêm lại: ``Tôi có cảm giác\dots\ những ngày bình yên này sẽ không kéo dài lâu đâu.''

Một thoáng im lặng.

``\dots Em biết rồi ạ.''

Tôi bước ra hành lang. Cửa phòng giáo viên khép lại phía sau, khe cửa hắt bóng hình tôi xuống nền gạch màu nâu sậm. Tôi đứng đó một lúc, tay vẫn còn đặt trên balo, đầu ong ong với những câu mà cô giáo chủ nhiệm chúng tôi vừa nói. Những tiếng trẻ con khóc. Những mảnh ký ức chia đôi. Đứa trẻ tan thành bụi sáng. Ánh mắt của Yuko-sensei.

Tất cả xoáy lại thành một vòng xoắn.

``\textit{Mình không thể bỏ qua chuyện này.
  Những đứa trẻ đó\dots\ chắc chắn còn điều gì đó muốn nói.}''

Tôi khẽ thở dài, rồi bước xuống cầu thang, hòa vào nhịp sống học đường bình thường đang diễn ra phía dưới, những âm thanh cười nói, tiếng giày lẹp kẹp, tiếng bóng nảy ngoài sân.

Nhưng trong lòng tôi đã rõ: từ những ngày tới, tôi sẽ tiếp tục điều tra. Về những linh hồn trẻ con ấy. Và về những vết thương âm ỉ mà thị trấn Aogami này vẫn còn tồn tại.

``\eng{Cậu đã quyết định rồi chứ?}'' giọng Game vang lên từ cổ tay, nhỏ như tiếng gió lướt qua. Như thể ông ta vừa đọc được suy nghĩ của tôi.

Tôi siết nhẹ đồng hồ. ``\eng{Ừ. Không thể lùi được nữa.}''

``\eng{Ngay cả khi cô giáo kia đã cảnh báo?}''

``\eng{Chính vì cô ấy đã cảnh báo, nên ta càng không thể ngồi yên.}'' Tôi bước chậm, né tránh một nhóm học sinh đang cười đùa. ``\eng{Nếu Yuko-sensei biết quá nhiều, không chắc chuyện gì sẽ xảy ra với cô ấy nữa.}''

``\eng{Cậu lo cho cô ta?}'' ông ta cười nhạt. ``\eng{Hay lo cho chính mình?}''

Tôi dừng lại bậc thang cuối, mắt dõi theo bóng chiều in trên sân. ``\eng{Cả hai.}''

Tiếng Game như tan trong không khí: ``\eng{Vậy thì đừng chết sớm, tổng quản. Cậu đang đánh cược cả mạng sống mình vào chuyện này đấy.}''

Tôi không đáp lại. Chỉ để lại sau lưng dãy hành lang dài những bước chắc nịch, hòa vào dòng người tan học, nhưng lòng đã rẽ sang một con đường khác, nơi những linh hồn lạc lối đang chờ đợi một lời giải đáp.
\newpage

\trva

Dưới đây là danh sách các học sinh của lớp 1-C và các câu lạc bộ mà họ gia nhập:

\begin{center}
  \begin{tabular}{|c|c|c|c|}
    \hline
    \textbf{STT} & \textbf{Tên đầy đủ} & \textbf{ID} & \textbf{Câu lạc bộ}       \\ \hline
    1            & Inukai Akane        & H4N045S004  & Nhiếp ảnh                 \\ \hline
    2            & Mizuno Aku          & H4N045S439  & Huyền bí                  \\ \hline
    3            & Utachi Azuki        & H4N045S312  & Âm nhạc                   \\ \hline
    4            & Onidoro Ayame       & H4N045S315  & Hội học sinh              \\ \hline
    5            & Shirayuki Inari     & H4N045S012  & Hội học sinh              \\ \hline
    6            & Mizuta Uzuki        & H4N045S005  & Nhiếp ảnh                 \\ \hline
    7            & Hikawa Kaigou       & H4N045S121  & Bóng rổ                   \\ \hline
    8            & Ukai Kaoru          & H4N045S008  & Phát thanh (khách mời)    \\ \hline
    9            & Tokiwa Kana         & H4N045S317  & Phát thanh                \\ \hline
    10           & Amano Kanade        & H4N045S330  & Phát thanh (khách mời)    \\ \hline
    11           & Hikawa Kuon         & H4N045S122  & "Về nhà" (không tham gia) \\ \hline
    12           & Akaba Koishin       & H4N045S336  & Văn học                   \\ \hline
    13           & Ibara Saki          & H4N045S020  & Hội chăm sóc cây cảnh     \\ \hline
    14           & Shinomiya Sakura    & H4N045S310  & Huyền bí                  \\ \hline
    15           & Kisaragi Saya       & H4N045S029  & Thiên văn                 \\ \hline
    16           & Shitou Shion        & H4N045S450  & Huyền bí                  \\ \hline
    17           & Fukami Shiina       & H4N045S311  & Bóng rổ                   \\ \hline
    18           & Misora Shino        & H4N045S016  & Hội học sinh              \\ \hline
    19           & Hikawa Suzuki       & H4N045S123  & Văn học                   \\ \hline
    20           & Karamomo Suzune     & H4N045S026  & Thiên văn                 \\ \hline
    21           & Yukihara Touka      & H4N045S027  & Hội chăm sóc cây cảnh     \\ \hline
    22           & Furukawa Nanase     & H4N045S323  & Hội học sinh              \\ \hline
    23           & Natsuno Hanabi      & H4N045S019  & Thiết kế thời trang       \\ \hline
    24           & Shiratori Hotaru    & H4N045S313  & Mỹ thuật                  \\ \hline
    25           & Omori Honoka        & H4N045S309  & Âm nhạc                   \\ \hline
    26           & Akagane Mari        & H4N045S322  & Mỹ thuật                  \\ \hline
    27           & Kamishiro Miku      & H4N045S007  & Thiết kế thời trang       \\ \hline
    28           & Himekawa Mitsuki    & H4N045S321  & Thiết kế thời trang       \\ \hline
    29           & Shiromi Yae         & H4N045S314  & Bóng rổ                   \\ \hline
    30           & Nekomiya Yuka       & H4N045S002  & Huyền bí                  \\ \hline
    31           & Kaizuka Yuki        & H4N045S338  & Âm nhạc                   \\ \hline
    32           & Hitsujikado Wata    & H4N045S226  & Hội học sinh              \\ \hline
  \end{tabular}
\end{center}

\end{document}